% =============================================================================
%  Monografia (Trabalho de Conclusão de Curso) — Robofael
%  Formatação segundo a ABNT: NBR 14724 (apresentação), NBR 6024 (numeração
%  progressiva), NBR 6027 (sumário), NBR 6028 (resumo), NBR 6023 (referências)
%  e NBR 10520 (citações, sistema autor-data).
%  Compilar com XeLaTeX:  make monografia   (ou: latexmk -xelatex monografia.tex)
% =============================================================================
\documentclass[12pt,a4paper,oneside]{report}

% ---------- Dados do trabalho (edite aqui) ----------
\newcommand{\autor}{Rafael Lourenço}
\newcommand{\titulo}{ROBOFAEL}
\newcommand{\subtitulo}{assistente pessoal com avatar animado e interação por voz, construído com arquitetura hexagonal e princípios SOLID}
\newcommand{\instituicao}{[NOME DA INSTITUIÇÃO DE ENSINO]}
\newcommand{\unidade}{[NOME DA FACULDADE OU DEPARTAMENTO]}
\newcommand{\curso}{[NOME DO CURSO]}
\newcommand{\grau}{[bacharel em NOME DO CURSO]}
\newcommand{\orientacao}{[Nome completo do(a) orientador(a)]}
\newcommand{\local}{[Cidade]}
\newcommand{\ano}{2026}

% ---------- Página: A4, margens 3 cm (sup./esq.) e 2 cm (inf./dir.) ----------
\usepackage[a4paper,top=3cm,left=3cm,bottom=2cm,right=2cm,headheight=15pt,headsep=0.5cm]{geometry}

% ---------- Idioma e fontes (Liberation Serif = métrica da Times New Roman) ----------
\usepackage{fontspec}
% Sem padrões de hifenização do português nesta distribuição TeX: em vez de hifenizar
% com as regras do inglês, a hifenização é desligada (texto justificado com folga extra).
\hyphenpenalty=10000
\exhyphenpenalty=10000
\newenvironment{idiomaestrangeiro}{}{}
\setmainfont{Liberation Serif}
\setmonofont{Liberation Mono}[Scale=0.88]

% ---------- Espaçamento 1,5, recuo de parágrafo e justificação ----------
\usepackage{setspace}
\onehalfspacing
\usepackage{indentfirst}
\setlength{\parindent}{1.25cm}
\setlength{\parskip}{0pt}
\setlength{\emergencystretch}{3em}
\sloppy
\usepackage{microtype}

% ---------- Tabelas, quadros e listas ----------
\usepackage{array}
\usepackage{ltablex}
\keepXColumns
\usepackage{booktabs}
\usepackage{enumitem}
\setlist{nosep,leftmargin=1.9cm}
\renewcommand{\contentsname}{Sumário}
\renewcommand{\listtablename}{Lista de tabelas}
\newcommand{\fonte}[1]{\par\vspace{2pt}{\noindent\small\singlespacing Fonte: #1\par}\vspace{12pt}}
\newcolumntype{L}{>{\raggedright\arraybackslash}X}
\newcolumntype{C}[1]{>{\centering\arraybackslash}p{#1}}
\newcolumntype{P}[1]{>{\raggedright\arraybackslash}p{#1}}

% ---------- Títulos: numeração progressiva (NBR 6024) ----------
\usepackage{titlesec}
\titleformat{\chapter}[hang]{\normalfont\bfseries\normalsize}{\thechapter}{0.6em}{\MakeUppercase}
\titlespacing*{\chapter}{0pt}{0pt}{18pt}
% Títulos sem indicativo numérico (sumário, listas, referências...): centralizados
\titleformat{name=\chapter,numberless}[block]{\normalfont\bfseries\normalsize\filcenter}{}{0pt}{\MakeUppercase}
\titlespacing*{name=\chapter,numberless}{0pt}{0pt}{18pt}
\titleformat{\section}[hang]{\normalfont\normalsize}{\thesection}{0.6em}{\MakeUppercase}
\titlespacing*{\section}{0pt}{18pt}{12pt}
\titleformat{\subsection}[hang]{\normalfont\bfseries\normalsize}{\thesubsection}{0.6em}{}
\titlespacing*{\subsection}{0pt}{18pt}{12pt}
\titleformat{\subsubsection}[hang]{\normalfont\normalsize}{\thesubsubsection}{0.6em}{}
\titlespacing*{\subsubsection}{0pt}{12pt}{12pt}
\setcounter{secnumdepth}{3}
\setcounter{tocdepth}{3}

% ---------- Paginação: algarismos arábicos no canto superior direito ----------
\usepackage{fancyhdr}
\fancypagestyle{textual}{\fancyhf{}\fancyhead[R]{\small\thepage}\renewcommand{\headrulewidth}{0pt}}
\fancypagestyle{plain}{\fancyhf{}\renewcommand{\headrulewidth}{0pt}}

% ---------- Sumário (NBR 6027) e listas ----------
\usepackage[titles]{tocloft}  % títulos das listas via \chapter* (centralizados, maiúsculas)
\newlistof[chapter]{quadro}{loq}{Lista de quadros}
\renewcommand{\thequadro}{\arabic{quadro}}
\counterwithout{quadro}{chapter}
% Legendas acima do quadro/tabela (NBR 14724): "Quadro 1 – Título"
\newcommand{\legendaquadro}[1]{\refstepcounter{quadro}\par\nopagebreak
  {\centering\small\singlespacing Quadro \thequadro{} – #1\par}\nopagebreak\vspace{4pt}%
  \addcontentsline{loq}{quadro}{\protect\numberline{\thequadro}#1}}
% Listas sem o espaço extra que o LaTeX insere a cada capítulo
\makeatletter
\newcommand{\listadequadros}{\begingroup\let\addvspace\@gobble\listofquadro\endgroup}
\newcommand{\listadetabelas}{\begingroup\let\addvspace\@gobble\listoftables\endgroup}
\makeatother
% Contador próprio: o longtable também incrementa o contador "table" a cada tabela
\newcounter{tabela}
\newcommand{\legendatabela}[1]{\refstepcounter{tabela}\par\nopagebreak
  {\centering\small\singlespacing Tabela \thetabela{} – #1\par}\nopagebreak\vspace{4pt}%
  \addcontentsline{lot}{table}{\protect\numberline{\thetabela}#1}}
\renewcommand{\cftchapfont}{\bfseries}
\renewcommand{\cftchappagefont}{\bfseries}
\renewcommand{\cftchapleader}{\cftdotfill{\cftdotsep}}
\renewcommand{\cftchappresnum}{}
\setlength{\cftchapindent}{0pt}
\setlength{\cftsecindent}{0pt}
\setlength{\cftsubsecindent}{0pt}
\setlength{\cftsubsubsecindent}{0pt}
\setlength{\cftchapnumwidth}{2.2em}
\setlength{\cftsecnumwidth}{2.2em}
\setlength{\cftsubsecnumwidth}{3em}
\setlength{\cftsubsubsecnumwidth}{3.8em}
\setlength{\cftbeforechapskip}{4pt}
\renewcommand{\cftsubsecfont}{\bfseries}
\renewcommand{\cftsubsecpagefont}{\bfseries}
\renewcommand{\cfttabpresnum}{Tabela~}
\renewcommand{\cftquadropresnum}{Quadro~}
\setlength{\cfttabnumwidth}{5em}
\setlength{\cftquadronumwidth}{6.6em}
\renewcommand{\cfttabaftersnum}{~–}
\renewcommand{\cftquadroaftersnum}{~–}

% ---------- Hiperligações e utilitários ----------
\usepackage{etoolbox}
\usepackage{xcolor}
\PassOptionsToPackage{obeyspaces,spaces}{url}  % \cod preserva espaços (ex.: make verificar)
\usepackage[hidelinks,bookmarksnumbered]{hyperref}
\usepackage{url}
\urlstyle{same}
% Identificadores de código em fonte monoespaçada, com quebra de linha em _ / . -
\DeclareUrlCommand\cod{\urlstyle{tt}}
\renewcommand{\UrlBreaks}{\do\_\do\/\do\.\do\-}
\newcommand{\en}[1]{\textit{#1}}                  % termos em língua estrangeira

% ---------- Citações autor-data (NBR 10520:2023): (Autor, ano) e Autor (ano) ----------
\makeatletter
\newcommand{\defref}[3]{\expandafter\def\csname ref@#1@a\endcsname{#2}\expandafter\def\csname ref@#1@y\endcsname{#3}}
\newcommand{\refcampo}[2]{\ifcsname ref@#1@#2\endcsname\csname ref@#1@#2\endcsname
  \else\textbf{[?#1]}\PackageWarning{monografia}{Referencia indefinida: #1}\fi}
\newcommand{\@citeum}[1]{\@citesep\refcampo{#1}{a}, \refcampo{#1}{y}\gdef\@citesep{; }}
\AtBeginDocument{%
  \renewcommand{\cite}[1]{(\gdef\@citesep{}\forcsvlist{\@citeum}{#1})}%
  \providecommand{\textcite}{}\renewcommand{\textcite}[1]{\refcampo{#1}{a} (\refcampo{#1}{y})}}
\makeatother
% Chaves de citação: \defref{chave}{Autor como aparece na citação}{ano}
\defref{abnt6023}{Associação Brasileira de Normas Técnicas}{2018}
\defref{abnt10520}{Associação Brasileira de Normas Técnicas}{2023}
\defref{abnt14724}{Associação Brasileira de Normas Técnicas}{2011}
\defref{alphacephei}{Alpha Cephei}{2026}
\defref{anthropic}{Anthropic}{2026}
\defref{beck2003}{Beck}{2003}
\defref{cockburn2005}{Cockburn}{2005}
\defref{cohn2009}{Cohn}{2009}
\defref{debian}{Debian}{2026}
\defref{evans2004}{Evans}{2004}
\defref{ford2017}{Ford; Parsons; Kua}{2017}
\defref{fowler2004}{Fowler}{2004}
\defref{fowler2018}{Fowler}{2018}
\defref{freedesktop}{Freedesktop.org}{2020}
\defref{gamma1994}{Gamma \textit{et al.}}{1994}
\defref{harris2020}{Harris \textit{et al.}}{2020}
\defref{hunt1999}{Hunt; Thomas}{1999}
\defref{iso25010}{International Organization for Standardization}{2011}
\defref{jurafsky2009}{Jurafsky; Martin}{2009}
\defref{larman2003}{Larman; Basili}{2003}
\defref{liskov1994}{Liskov; Wing}{1994}
\defref{martin2003}{Martin}{2003}
\defref{martin2009}{Martin}{2009}
\defref{martin2019}{Martin}{2019}
\defref{mccabe1976}{McCabe}{1976}
\defref{meszaros2007}{Meszaros}{2007}
\defref{meyer1988}{Meyer}{1988}
\defref{nielsen1993}{Nielsen}{1993}
\defref{nist2015}{National Institute of Standards and Technology}{2015}
\defref{ollama}{Ollama}{2026}
\defref{pep544}{Levkivskyi; Lehtosalo; Langa}{2017}
\defref{pipewire}{PipeWire}{2026}
\defref{povey2011}{Povey \textit{et al.}}{2011}
\defref{pressman2016}{Pressman; Maxim}{2016}
\defref{semver}{Preston-Werner}{2013}
\defref{sohn1999}{Sohn; Kim; Sung}{1999}
\defref{sommerville2018}{Sommerville}{2018}
\defref{sqlite}{SQLite}{2026}
\defref{vaswani2017}{Vaswani \textit{et al.}}{2017}


% Sumário com a mesma grafia dos títulos (NBR 6027): capítulos e seções em maiúsculas.
% O texto da entrada é convertido ao ser gravado; o número da seção (\numberline) é preservado.
\makeatletter
\newcommand{\tocmaiusc}[1]{\@tocmaiusc#1\@nil}
\def\@tocmaiusc#1#2\@nil{\ifx#1\numberline\@tocnum#2\@nil\else\MakeUppercase{#1#2}\fi}
\def\@tocnum#1#2\@nil{\numberline{#1}\MakeUppercase{#2}}
\AtBeginDocument{%
  \let\addcontentslineoriginal\addcontentsline
  \renewcommand{\addcontentsline}[3]{%
    \ifboolexpr{test {\ifstrequal{#1}{toc}} and (test {\ifstrequal{#2}{chapter}} or test {\ifstrequal{#2}{section}})}%
      {\addcontentslineoriginal{#1}{#2}{\texorpdfstring{\protect\tocmaiusc{#3}}{#3}}}%
      {\addcontentslineoriginal{#1}{#2}{#3}}}}
\makeatother

% Título sem indicativo numérico, centralizado (elementos pré e pós-textuais)
\newcommand{\titulosemnumero}[1]{\clearpage\begin{center}\textbf{\MakeUppercase{#1}}\end{center}\vspace{12pt}}
\newcommand{\titulopostextual}[1]{\clearpage\phantomsection\addcontentsline{toc}{chapter}{\MakeUppercase{#1}}%
  \begin{center}\textbf{\MakeUppercase{#1}}\end{center}\vspace{12pt}}

\begin{document}
\pagestyle{empty}

% =============================================================================
%  Elementos pré-textuais (NBR 14724): capa, folha de rosto, folha de aprovação,
%  resumo, abstract, listas e sumário. Não são paginados, mas a contagem começa
%  na folha de rosto.
% =============================================================================

% ---------- Capa (não é contada) ----------
\begin{titlepage}
\begin{center}
\singlespacing
\textbf{\MakeUppercase{\instituicao}}\\
\textbf{\MakeUppercase{\unidade}}\\
\textbf{\MakeUppercase{\curso}}

\vspace*{3.2cm}
\textbf{\MakeUppercase{\autor}}

\vfill
\textbf{\titulo:} \subtitulo

\vfill
\local\\
\ano
\end{center}
\end{titlepage}

% ---------- Folha de rosto (página 1, não numerada) ----------
\setcounter{page}{1}
\begin{center}
\singlespacing
\textbf{\MakeUppercase{\autor}}

\vspace*{5cm}
\textbf{\titulo:} \subtitulo
\end{center}

\vspace{2.5cm}
\noindent\hspace{0.45\textwidth}\begin{minipage}{0.55\textwidth}
\singlespacing\small
Trabalho de Conclusão de Curso apresentado ao \curso{} da \instituicao, como requisito parcial
para a obtenção do título de \grau.

\vspace{12pt}
Orientador(a): \orientacao.
\end{minipage}

\vfill
\begin{center}
\local\\
\ano
\end{center}

% ---------- Folha de aprovação ----------
\clearpage
\begin{center}
\singlespacing
\textbf{\MakeUppercase{\autor}}

\vspace{1.5cm}
\textbf{\titulo:} \subtitulo
\end{center}

\vspace{1cm}
\noindent\hspace{0.45\textwidth}\begin{minipage}{0.55\textwidth}
\singlespacing\small
Trabalho de Conclusão de Curso apresentado ao \curso{} da \instituicao, como requisito parcial
para a obtenção do título de \grau.
\end{minipage}

\vspace{1cm}
\noindent Aprovado em: \underline{\hspace{1cm}} / \underline{\hspace{1cm}} / \underline{\hspace{1.6cm}}.

\vspace{1cm}
\begin{center}
\textbf{BANCA EXAMINADORA}

\vspace{1.6cm}
\rule{10cm}{0.4pt}\\
\orientacao{} (Orientador(a))\\
\instituicao

\vspace{1.4cm}
\rule{10cm}{0.4pt}\\
{[Nome completo do(a) avaliador(a)]}\\
{[Instituição]}

\vspace{1.4cm}
\rule{10cm}{0.4pt}\\
{[Nome completo do(a) avaliador(a)]}\\
{[Instituição]}
\end{center}

% ---------- Resumo na língua vernácula (NBR 6028) ----------
\titulosemnumero{Resumo}
\noindent
Este trabalho apresenta o projeto, a construção e a validação do Robofael, um assistente pessoal de
atividades diárias para computadores com Linux que conversa por voz e é representado por um avatar
animado gerado a partir de fotografias do próprio usuário. O sistema ouve comandos falados ou
digitados, responde com voz neural ajustada ao tom e ao ritmo da voz do usuário, gerencia tarefas e
lembretes, mantém uma memória de longo prazo auditável capaz de responder o que foi feito em um dia ou
em uma semana, abre programas instalados e funciona em português e em inglês. A ativação sem as mãos
ocorre pela frase “ok Fael”, reconhecida localmente, de modo que o áudio do ambiente não deixa o
computador. Quando disponível, um modelo de linguagem de grande porte (Claude) interpreta pedidos em
linguagem natural e opera a agenda por meio de ferramentas; sem ele, um modelo local servido pelo Ollama
conversa com ferramentas restritas e, sem nenhum dos dois, um interpretador de intenções local mantém as
funções essenciais. O desenvolvimento seguiu um processo iterativo e incremental,
conduzido por pedidos sucessivos e apoiado por um agente de programação baseado em inteligência
artificial, em vinte incrementos registrados em versionamento semântico. A solução adota arquitetura
hexagonal, pipeline de fala em fluxo contínuo, máquina de estados explícita, barramento de eventos
para isolar a interface gráfica e persistência em SQLite com histórico somente de acréscimo protegido
por cadeia de resumos criptográficos SHA-256. Ao final, o código foi revisado de ponta a ponta à luz
dos princípios SOLID e das práticas de código limpo: responsabilidades foram separadas, despachos
condicionais foram convertidos em tabelas extensíveis e as dependências passaram a apontar apenas
para contratos. As regras de arquitetura tornaram-se verificáveis por testes automatizados. Um
incremento guiado por medições reduziu o consumo de processamento em repouso de 29\% para 5\% de um
núcleo, levou a janela a aparecer em 0,6 s e eliminou pausas artificiais entre as frases faladas. A
validação contou com 643 testes automatizados, cobertura de 87\% das linhas e ramificações,
verificação estática de tipos e de estilo, complexidade ciclomática máxima igual a 10 por função e um
instalador para Ubuntu testado em ambiente isolado. Os resultados indicam que é viável construir um
assistente de voz responsivo, privado e de fácil manutenção com componentes livres e serviços de
linguagem, e que a verificação automática das regras de projeto preserva a qualidade da arquitetura
ao longo da evolução do software.

\vspace{12pt}
\noindent\textbf{Palavras-chave:} assistente virtual; reconhecimento de fala; arquitetura hexagonal;
princípios SOLID; código limpo.

% ---------- Resumo em língua estrangeira ----------
\titulosemnumero{Abstract}
\begin{idiomaestrangeiro}
\noindent
This work presents the design, construction and validation of Robofael, a personal daily-activity
assistant for Linux computers that talks by voice and is represented by an animated avatar generated
from photographs of its own user. The system listens to spoken or typed commands, answers with a neural
voice tuned to the pitch and pace of the user's voice, manages tasks and reminders, keeps an auditable
long-term memory able to answer what was done on a given day or week, opens installed programs and works
in Portuguese and English. Hands-free activation happens through the phrase “ok Fael”, recognized
locally, so ambient audio never leaves the computer. When available, a large language model (Claude)
interprets natural-language requests and operates the agenda through tools; without it, a local model
served by Ollama talks with restricted tools and, without either, a local intent interpreter keeps the
essential functions. Development followed an iterative and incremental process,
driven by successive requests and supported by an artificial-intelligence coding agent, in twenty
increments tracked with semantic versioning. The solution adopts a hexagonal architecture, a streaming
speech pipeline, an explicit state machine, an event bus that isolates the graphical interface, and
SQLite persistence with an append-only history protected by a SHA-256 hash chain. Finally, the code was
reviewed end to end in the light of the SOLID principles and clean-code practices: responsibilities were
separated, conditional dispatch was replaced with extensible tables, dependencies were made to point only
to contracts, and architecture rules became verifiable through automated tests. A measurement-driven
increment cut idle processor usage from 29\% to 5\% of one core, made the window appear in 0.6 s and
removed artificial pauses between spoken sentences. Validation relied on
643 automated tests, 87\% line and branch coverage, static type and style checking, a maximum cyclomatic
complexity of 10 per function, and an Ubuntu installer tested in an isolated environment. The results
indicate that it is feasible to build a responsive, private and maintainable voice assistant from free
components and language services, and that automatically checking design rules preserves architectural
quality as the software evolves.

\vspace{12pt}
\noindent\textbf{Keywords:} virtual assistant; speech recognition; hexagonal architecture; SOLID
principles; clean code.
\end{idiomaestrangeiro}

% ---------- Listas ----------
\clearpage
\listadequadros
\clearpage
\listadetabelas

% ---------- Lista de abreviaturas e siglas ----------
\titulosemnumero{Lista de abreviaturas e siglas}
\begin{singlespace}
\noindent
\begin{tabular}{@{}p{2.6cm}p{12cm}@{}}
ABNT   & Associação Brasileira de Normas Técnicas \\
AGC    & \en{Automatic Gain Control} (controle automático de ganho) \\
API    & \en{Application Programming Interface} (interface de programação de aplicações) \\
CI     & \en{Continuous Integration} (integração contínua) \\
DIP    & \en{Dependency Inversion Principle} (princípio da inversão de dependência) \\
DSP    & \en{Digital Signal Processing} (processamento digital de sinais) \\
FPS    & Quadros por segundo \\
GIL    & \en{Global Interpreter Lock} do interpretador CPython \\
GUI    & \en{Graphical User Interface} (interface gráfica do usuário) \\
ISP    & \en{Interface Segregation Principle} (princípio da segregação de interfaces) \\
JSON   & \en{JavaScript Object Notation} \\
LLM    & \en{Large Language Model} (modelo de linguagem de grande porte) \\
LRU    & \en{Least Recently Used} (menos recentemente usado) \\
LSP    & \en{Liskov Substitution Principle} (princípio da substituição de Liskov) \\
OCP    & \en{Open-Closed Principle} (princípio aberto/fechado) \\
PCM    & \en{Pulse-Code Modulation} (modulação por código de pulso) \\
RF     & Requisito funcional \\
RNF    & Requisito não funcional \\
SHA    & \en{Secure Hash Algorithm} \\
SOLID  & Acrônimo dos cinco princípios de projeto orientado a objetos (SRP, OCP, LSP, ISP, DIP) \\
SQL    & \en{Structured Query Language} \\
SRP    & \en{Single Responsibility Principle} (princípio da responsabilidade única) \\
STT    & \en{Speech-to-Text} (reconhecimento de fala) \\
TTS    & \en{Text-to-Speech} (síntese de fala) \\
UI     & \en{User Interface} (interface do usuário) \\
VAD    & \en{Voice Activity Detection} (detecção de atividade de voz) \\
XDG    & Especificações de área de trabalho do projeto freedesktop.org \\
\end{tabular}
\end{singlespace}

% ---------- Sumário ----------
\clearpage
\tableofcontents


% ---------- Parte textual: paginação visível a partir daqui ----------
\clearpage
\pagestyle{textual}
\fancypagestyle{plain}{\fancyhf{}\fancyhead[R]{\small\thepage}\renewcommand{\headrulewidth}{0pt}}

\chapter{Introdução}

Assistentes pessoais por voz tornaram-se parte do cotidiano em telefones celulares, caixas de som
e automóveis. Em computadores pessoais, porém, sobretudo nos que utilizam o sistema operacional
Linux, a oferta de assistentes integrados ao ambiente de trabalho ainda é limitada. As soluções
comerciais costumam depender de serviços em nuvem permanentemente à escuta, oferecem pouca
personalização da aparência e da voz e não guardam um histórico verificável das atividades
realizadas pelo usuário. Ao mesmo tempo, o surgimento de modelos de linguagem de grande porte
(\en{Large Language Models}, LLM), capazes de interpretar pedidos em linguagem natural e de acionar
funções externas, reduziu drasticamente o esforço necessário para construir um assistente que
compreenda frases livres, e não apenas comandos rígidos \cite{vaswani2017}.

Construir um assistente de voz útil, entretanto, envolve muito mais do que conectar um reconhecedor
de fala a um modelo de linguagem. É preciso capturar o áudio com qualidade, decidir quando o usuário
terminou de falar, sintetizar a resposta sem longas esperas, sincronizar a boca de um avatar com a
voz, persistir dados de forma confiável e manter a interface gráfica fluida enquanto diversas tarefas
concorrentes acontecem. Esse conjunto de preocupações, distribuído entre várias tecnologias, torna o
software suscetível a acoplamento excessivo e à degradação de sua estrutura à medida que novas
funções são acrescentadas, fenômeno descrito por \textcite{martin2019} como o principal fator do custo
crescente de manutenção.

Este trabalho descreve a concepção, a construção incremental e a validação do \textbf{Robofael}, um
assistente pessoal de atividades diárias que roda localmente em um computador com Linux. O Robofael
é representado por um avatar animado criado a partir de fotografias do próprio usuário, ouve e
responde por voz, organiza tarefas e lembretes, responde perguntas sobre o que foi feito em
determinado período, abre programas instalados e conversa em português ou em inglês. Além da
descrição funcional, o trabalho dá ênfase à qualidade interna do software: a solução foi organizada
segundo a arquitetura hexagonal \cite{cockburn2005} e, ao final, revisada de ponta a ponta à luz dos
princípios SOLID \cite{martin2003} e das práticas de código limpo \cite{martin2009}, com regras de
projeto verificadas automaticamente por testes.

\section{Problema de pesquisa}

O problema que orienta este trabalho pode ser enunciado da seguinte forma: \emph{como construir um
assistente pessoal por voz para computadores com Linux que seja responsivo, preserve a privacidade do
usuário, mantenha um histórico confiável de suas atividades e, ao mesmo tempo, conserve uma estrutura
interna fácil de entender, testar e evoluir?}

\section{Objetivos}

\subsection{Objetivo geral}

Desenvolver e validar um assistente pessoal de atividades diárias, com avatar animado e interação por
voz, cuja arquitetura siga os princípios SOLID e as práticas de código limpo, com comportamento e
regras de projeto verificados por testes automatizados.

\subsection{Objetivos específicos}

\begin{enumerate}[label=\alph*)]
  \item gerar, a partir de fotografias do usuário, um avatar em formato de busto com movimentação da
        boca sincronizada à voz sintetizada;
  \item implementar o ciclo de conversa por voz (ouvir, compreender e responder) com latência
        percebida próxima à da primeira frase da resposta;
  \item permitir a ativação sem as mãos por uma frase de chamada reconhecida localmente;
  \item oferecer gestão de tarefas e lembretes e uma memória de longo prazo auditável, capaz de
        responder o que foi feito em um dia, em uma semana ou em outro período;
  \item integrar o assistente ao ambiente Linux, permitindo abrir programas instalados por comando de
        voz;
  \item suportar os idiomas português do Brasil e inglês dos Estados Unidos, com troca em tempo de
        execução;
  \item disponibilizar um instalador padrão para o Ubuntu;
  \item revisar a solução segundo os princípios SOLID e as práticas de código limpo, tornando as
        regras de arquitetura verificáveis por testes automatizados;
  \item documentar os requisitos funcionais e não funcionais e todo o processo de construção.
\end{enumerate}

\section{Justificativa}

A relevância do trabalho manifesta-se em três dimensões. Na dimensão do usuário, um assistente que
roda no próprio computador, escuta passivamente apenas de forma local e guarda os dados em um arquivo
sob controle do usuário atende a uma preocupação crescente com privacidade. A memória de longo prazo
auditável, por sua vez, transforma o assistente em uma ferramenta de acompanhamento de produtividade:
ao perguntar “o que eu fiz na semana passada?”, o usuário obtém uma resposta baseada em registros
cuja integridade pode ser verificada.

Na dimensão técnica, o projeto reúne, em um único sistema de porte moderado, problemas clássicos de
engenharia de software: concorrência entre captura de áudio, geração de texto, síntese de voz e
animação; integração com serviços externos sujeitos a falhas; persistência transacional; e
internacionalização. Isso o torna um estudo de caso adequado para avaliar como princípios de projeto
orientado a objetos se comportam em um contexto real.

Na dimensão acadêmica, o trabalho documenta um processo de desenvolvimento conduzido com o apoio de
um agente de programação baseado em inteligência artificial e mostra que a qualidade estrutural, nesse
cenário, pode ser preservada por meio de verificações automáticas, como testes de arquitetura, limites
de complexidade e verificação estática de tipos, em vez de depender exclusivamente de disciplina
manual.

\section{Delimitação do escopo}

O sistema destina-se a um único usuário por conta do sistema operacional e foi construído e testado
no Ubuntu 24.04 com Python 3.12, embora o código seja compatível com Python 3.10 ou superior. O
reconhecimento de fala para comandos usa o serviço de reconhecimento do Google acessado por meio da
biblioteca SpeechRecognition, e a síntese de voz utiliza vozes neurais do serviço edge-tts, com o gTTS
como reserva; ambos exigem conexão com a internet. A interpretação em linguagem natural usa a API do
Claude ou, sem ela, um modelo local servido pelo Ollama, mais lento quando executado no processador; o
modo offline cobre um conjunto definido de comandos. O treinamento de
modelos acústicos próprios, a clonagem completa do timbre de voz e o suporte a outros sistemas
operacionais estão fora do escopo deste trabalho.

\section{Organização do trabalho}

Além desta introdução, o trabalho está organizado em nove capítulos. O Capítulo 2 apresenta a
fundamentação teórica sobre assistentes de voz, processamento de fala, modelos de linguagem,
arquitetura de software, princípios SOLID e código limpo. O Capítulo 3 descreve a metodologia de
desenvolvimento adotada. O Capítulo 4 especifica os requisitos funcionais e não funcionais. O
Capítulo 5 detalha a arquitetura e os fluxos de execução. O Capítulo 6 relata, passo a passo, a
construção incremental do sistema. O Capítulo 7 apresenta a revisão do código à luz dos princípios
SOLID e do código limpo. O Capítulo 8 trata da estratégia de testes e da validação. O Capítulo 9
discute os resultados obtidos, e o Capítulo 10 apresenta as conclusões e os trabalhos futuros.

\chapter{Fundamentação teórica}

Este capítulo reúne os conceitos que sustentam as decisões tomadas no projeto. Inicialmente são
tratados os assistentes de voz e as tecnologias de fala; em seguida, os modelos de linguagem e o uso
de ferramentas; por fim, os fundamentos de engenharia de software que orientaram a organização do
código: arquitetura hexagonal, princípios SOLID, código limpo, métricas de complexidade, testes
automatizados e qualidade de software.

\section{Assistentes pessoais por voz}

Um assistente pessoal por voz é um sistema interativo que recebe pedidos falados, interpreta a
intenção do usuário, executa ações e devolve uma resposta, também falada. Segundo
\textcite{jurafsky2009}, sistemas de diálogo orientados a tarefas combinam tradicionalmente três
componentes: o reconhecimento automático de fala, que converte o sinal acústico em texto; a
compreensão da linguagem, que identifica a intenção e os parâmetros do pedido; e a geração e a síntese
da resposta. Entre esses componentes, um gerenciador de diálogo decide quando ouvir, quando responder
e como lidar com interrupções.

A qualidade percebida de um assistente depende fortemente da latência entre o fim da fala do usuário
e o início da resposta. \textcite{nielsen1993} associa tempos de resposta próximos de um segundo à
manutenção do fluxo de pensamento do usuário; acima de dez segundos, a atenção se perde. Em um
assistente de voz, a espera inclui a detecção do fim da fala, o reconhecimento, a geração do texto e
a síntese do áudio, o que motiva arquiteturas que processam essas etapas de forma sobreposta, em vez
de sequencial.

\section{Processamento de fala}

\subsection{Detecção de atividade de voz}

A detecção de atividade de voz (\en{Voice Activity Detection}, VAD) separa os trechos com fala dos
trechos com silêncio ou ruído. Abordagens clássicas comparam a energia do sinal com um limiar estimado a
partir do ruído de fundo, enquanto abordagens estatísticas modelam a distribuição do ruído e da fala
\cite{sohn1999}. No Robofael, a detecção é usada para decidir quando o usuário começou e terminou uma
frase. Para isso, calibra-se o limiar a partir do ruído medido na abertura do microfone e exige-se uma
pausa mínima para encerrar a captura. Ruídos curtos, como cliques, são descartados.

\subsection{Reconhecimento automático de fala}

O reconhecimento automático de fala (\en{Speech-to-Text}, STT) converte o áudio em texto. Sistemas
modernos combinam modelos acústicos e de linguagem treinados em grandes volumes de dados
\cite{jurafsky2009}. O projeto emprega dois reconhecedores com papéis complementares. O primeiro é o
Vosk, construído sobre o conjunto de ferramentas Kaldi \cite{povey2011}, que roda inteiramente no
computador e permite restringir o vocabulário a uma gramática fechada \cite{alphacephei}. O segundo é
o serviço de reconhecimento do Google, mais preciso para frases livres, usado apenas depois que o
usuário chama o assistente. A gramática restrita é especialmente útil para frases de ativação curtas,
que um reconhecedor de vocabulário aberto tende a confundir com palavras mais frequentes.

\subsection{Síntese de fala e prosódia}

A síntese de fala (\en{Text-to-Speech}, TTS) produz áudio a partir de texto. Vozes neurais atuais
permitem ajustar parâmetros prosódicos, como a frequência fundamental (o tom) e a velocidade de
elocução (o ritmo). A frequência fundamental de uma voz pode ser estimada pela autocorrelação do sinal
em janelas curtas, procurando o período que maximiza a semelhança do sinal consigo mesmo dentro da
faixa típica da voz humana, entre 60 Hz e 400 Hz \cite{jurafsky2009}. O Robofael usa essa estimativa
para comparar a voz gravada do usuário com a voz sintética e calcular o deslocamento de tom e de ritmo
que as aproxima.

\section{Modelos de linguagem de grande porte e uso de ferramentas}

Os modelos de linguagem de grande porte baseiam-se na arquitetura \en{Transformer}, cujo mecanismo de
atenção permite modelar dependências de longo alcance no texto \cite{vaswani2017}. Além de gerar texto,
esses modelos podem ser instruídos a solicitar a execução de funções externas, as chamadas ferramentas
(\en{tools}): o modelo devolve o nome da ferramenta e os argumentos em formato estruturado; a aplicação
executa a função e devolve o resultado; e o modelo continua a resposta. A API do Claude oferece esse
mecanismo com transmissão em fluxo contínuo (\en{streaming}) do texto e das entradas das ferramentas
\cite{anthropic}. Quando as entradas são transmitidas antes de estarem completas, cabe à aplicação
validá-las contra o esquema declarado antes de executar qualquer ação, porque uma entrada truncada não
pode ser tratada como válida.

Modelos menores, com alguns bilhões de parâmetros, podem ser executados no próprio computador por
servidores locais como o Ollama, que expõe uma API HTTP de conversa com suporte a ferramentas
\cite{ollama}. A execução local dispensa a conexão e não envia o texto a terceiros, mas, sem placa de
vídeo, as respostas são mais lentas, e modelos pequenos tendem a escolher ferramentas com menos
precisão. Neste trabalho, essa diferença foi tratada restringindo as ferramentas oferecidas ao modelo
local (Capítulo 6).

\section{Arquitetura de software}

\subsection{Arquitetura hexagonal}

A arquitetura hexagonal, também chamada de arquitetura de portas e adaptadores, foi proposta por
\textcite{cockburn2005} com o objetivo de isolar a lógica da aplicação dos mecanismos de entrada e
saída. A aplicação define \emph{portas}, isto é, interfaces que descrevem o que ela precisa do mundo
externo ou o que oferece a ele, e cada tecnologia concreta é conectada por meio de um
\emph{adaptador} que implementa a porta. Com isso, banco de dados, interface gráfica, serviços de
voz e modelos de linguagem tornam-se substituíveis, e a lógica central pode ser testada com dublês
simples. \textcite{martin2019} generaliza a ideia na chamada regra de dependência: as dependências de
código-fonte devem apontar sempre para dentro, das camadas de detalhe para as camadas de política.

\subsection{Domínio, aplicação e raiz de composição}

\textcite{evans2004} propõe que as regras de negócio sejam expressas em um modelo de domínio
independente de infraestrutura, com entidades, objetos de valor e serviços de domínio. Os casos de uso
que coordenam esse modelo formam a camada de aplicação. A ligação entre abstrações e implementações
concretas é feita por injeção de dependências \cite{fowler2004}, idealmente concentrada em um único
ponto do programa, a raiz de composição, que monta o grafo de objetos na inicialização.

\subsection{Padrões de projeto empregados}

Diversos padrões catalogados por \textcite{gamma1994} aparecem na solução: o \en{Observer}, na forma
de um barramento de eventos que desacopla produtores e consumidores; o \en{State}, em uma máquina de
estados explícita que controla o ciclo de conversa; o \en{Strategy} e a \en{Chain of Responsibility},
nas cadeias de interpretadores de intenções e de períodos; e o \en{Adapter}, nos componentes que
encapsulam bibliotecas externas.

\subsection{Funções de aptidão arquitetural}

\textcite{ford2017} propõem tratar características arquiteturais desejadas como funções de aptidão
(\en{fitness functions}): verificações objetivas e, preferencialmente, automatizadas, executadas
continuamente para indicar se a arquitetura ainda atende ao que se espera dela. Um exemplo é um teste
que inspeciona as dependências entre módulos e falha quando uma camada interna passa a depender de uma
externa. Essa ideia fundamenta os testes de arquitetura descritos no Capítulo 7.

\section{Princípios SOLID}

O acrônimo SOLID reúne cinco princípios de projeto orientado a objetos sistematizados por
\textcite{martin2003}:

\begin{enumerate}[label=\alph*)]
  \item \textbf{princípio da responsabilidade única} (\en{Single Responsibility Principle}, SRP): um
        módulo deve ter um, e apenas um, motivo para mudar; em outras palavras, deve responder a um único
        ator ou preocupação \cite{martin2019};
  \item \textbf{princípio aberto/fechado} (\en{Open-Closed Principle}, OCP): entidades de software
        devem estar abertas para extensão e fechadas para modificação, formulação atribuída a
        \textcite{meyer1988}; acrescentar um comportamento não deve exigir a alteração de código que já
        funciona;
  \item \textbf{princípio da substituição de Liskov} (\en{Liskov Substitution Principle}, LSP):
        objetos de um subtipo devem poder substituir objetos do tipo base sem alterar as propriedades
        desejáveis do programa \cite{liskov1994};
  \item \textbf{princípio da segregação de interfaces} (\en{Interface Segregation Principle}, ISP):
        clientes não devem ser forçados a depender de métodos que não usam, o que favorece interfaces
        pequenas e específicas;
  \item \textbf{princípio da inversão de dependência} (\en{Dependency Inversion Principle}, DIP):
        módulos de alto nível não devem depender de módulos de baixo nível; ambos devem depender de
        abstrações.
\end{enumerate}

Na linguagem Python, as abstrações exigidas pelo DIP e pelo ISP podem ser expressas por protocolos
(\cod{typing.Protocol}), que definem contratos verificados estruturalmente pelo verificador de tipos,
sem exigir herança explícita \cite{pep544}. Essa característica é explorada em todo o projeto.

\section{Código limpo e refatoração}

\textcite{martin2009} caracteriza o código limpo como aquele que é fácil de ler e de modificar por
outras pessoas: nomes que revelam a intenção, funções pequenas que fazem uma única coisa, poucos
argumentos, ausência de duplicação e de efeitos colaterais ocultos, e tratamento de erros que não
obscurece a lógica. O princípio da não repetição (\en{Don't Repeat Yourself}) de
\textcite{hunt1999} complementa essas práticas. A refatoração, segundo \textcite{fowler2018}, é a
alteração da estrutura interna do software sem mudar seu comportamento observável, feita em pequenos
passos e protegida por testes automatizados. Entre as técnicas aplicadas neste trabalho estão a
extração de classe, a extração de função, a substituição de condicionais por polimorfismo ou por
tabelas de despacho e a introdução de objetos-parâmetro.

\section{Métricas de complexidade}

A complexidade ciclomática, proposta por \textcite{mccabe1976}, mede o número de caminhos linearmente
independentes no grafo de controle de uma função. Ela corresponde, de forma aproximada, ao número de
decisões mais um. Valores altos indicam funções difíceis de entender e de testar por completo, e o
limite de 10 sugerido pelo próprio autor é amplamente adotado como teto prático. Neste trabalho, a
métrica foi usada tanto para localizar os pontos que precisavam de refatoração quanto como regra
permanente do analisador estático.

\section{Testes automatizados}

O desenvolvimento orientado a testes, descrito por \textcite{beck2003}, popularizou a prática de
escrever testes automatizados pequenos e rápidos que documentam o comportamento esperado.
\textcite{cohn2009} propõe a pirâmide de testes: muitos testes unitários na base, uma quantidade menor
de testes de integração e poucos testes de ponta a ponta, que são mais lentos e frágeis.
\textcite{meszaros2007} cataloga os dublês de teste, como \en{stubs}, \en{fakes} e \en{spies}, que
substituem colaboradores reais para isolar o comportamento sob teste. A combinação de arquitetura
hexagonal com dublês permite testar a lógica de conversa sem microfone, sem alto-falante e sem rede.

\section{Qualidade de software}

A norma ISO/IEC 25010 define um modelo de qualidade de produto com características como adequação
funcional, eficiência de desempenho, compatibilidade, usabilidade, confiabilidade, segurança,
manutenibilidade e portabilidade \cite{iso25010}. Esse modelo foi usado como referência para organizar
os requisitos não funcionais do Capítulo 4. Na engenharia de requisitos, \textcite{sommerville2018} e
\textcite{pressman2016} distinguem os requisitos funcionais, que descrevem o que o sistema deve fazer,
dos requisitos não funcionais, que restringem como ele o faz, distinção também adotada neste
trabalho.

\section{Integridade de dados e históricos auditáveis}

Um histórico é auditável quando permite reconstruir o que aconteceu e detectar alterações indevidas.
Duas técnicas complementares sustentam essa propriedade no projeto. A primeira é o armazenamento
somente de acréscimo: os registros nunca são alterados nem removidos, e restrições do próprio banco de
dados recusam essas operações. A segunda é o encadeamento de resumos criptográficos: cada registro
guarda o resumo SHA-256 \cite{nist2015} do anterior, de modo que qualquer modificação, remoção ou
inserção indevida quebra a cadeia a partir do ponto alterado. O banco de dados SQLite oferece
transações atômicas e gatilhos (\en{triggers}) adequados para implementar essas garantias em um único
arquivo local \cite{sqlite}.

\chapter{Metodologia}

Este capítulo descreve o tipo de pesquisa, o processo de desenvolvimento, o ambiente e as ferramentas
utilizadas e os critérios adotados para validar o sistema.

\section{Caracterização da pesquisa}

Quanto à natureza, trata-se de pesquisa aplicada, voltada à solução de um problema concreto por meio
da construção de um artefato de software. Quanto aos procedimentos, caracteriza-se como um estudo de
caso de desenvolvimento, no qual o próprio processo de construção, as decisões de projeto e os
problemas encontrados constituem objeto de análise. A abordagem é predominantemente qualitativa,
complementada por medidas quantitativas de qualidade interna, como o número de testes, a cobertura de
código, a complexidade ciclomática e o tamanho dos módulos, e de desempenho, como os tempos de
inicialização, de resposta e de interrupção.

\section{Processo de desenvolvimento}

Adotou-se um processo iterativo e incremental \cite{larman2003}. Cada incremento partiu de um pedido
do usuário, que atuou simultaneamente como cliente e como responsável pelas decisões de produto, e
terminou com o software executável, testado e documentado. Ao todo, foram realizados vinte
incrementos, descritos no Capítulo 6. Cada um deles seguiu as etapas abaixo:

\begin{enumerate}[label=\alph*)]
  \item \textbf{entendimento do pedido}: identificação do objetivo e das restrições; quando havia mais
        de uma solução razoável com consequências diferentes para o usuário, a escolha era submetida a
        ele, como na decisão entre o reconhecimento local e o remoto da frase de ativação ou no
        ajuste apenas do tom da voz;
  \item \textbf{investigação}: leitura do código afetado, medições, como volume do microfone,
        transcrições reais e tempos de execução, e experimentos curtos para confirmar hipóteses;
  \item \textbf{implementação}: alterações pequenas e reversíveis, preservando a arquitetura existente;
  \item \textbf{verificação}: execução da suíte de testes, da análise estática de estilo e de
        complexidade (ruff) e da verificação de tipos (mypy); acréscimo de testes para todo
        comportamento novo ou defeito corrigido;
  \item \textbf{validação com o usuário}: execução real do assistente, com análise de registros
        detalhados (opção \cod{-v}) para calibrar parâmetros como limiares de detecção e ganho do
        microfone;
  \item \textbf{documentação}: atualização do arquivo README, do registro de mudanças (CHANGELOG) e do
        número de versão segundo o versionamento semântico \cite{semver}.
\end{enumerate}

\section{Desenvolvimento assistido por agente de inteligência artificial}

A implementação foi realizada com o apoio do Claude Code, um agente de programação baseado em modelo de
linguagem que lê e edita arquivos, executa comandos e testes e interage com o usuário em linguagem
natural. O usuário definiu os objetivos e as prioridades, tomou as decisões de produto e validou o
comportamento em uso real; o agente propôs e implementou as soluções, executou as verificações e
relatou os resultados. Esse arranjo impôs uma disciplina adicional: como uma alteração pode ser
produzida rapidamente, a garantia de que ela não degrada o sistema precisa vir de verificações
automáticas. Por isso, testes, limites de complexidade, verificação de tipos e testes de arquitetura
foram tratados como parte do produto, e não como atividades opcionais.

\section{Ambiente e ferramentas}

O Quadro~\ref{qua:ferramentas} relaciona as tecnologias utilizadas e a finalidade de cada uma.

\begin{singlespace}
\legendaquadro{Tecnologias e ferramentas utilizadas}\label{qua:ferramentas}
\small
\begin{tabularx}{\textwidth}{|P{4.3cm}|L|}
\hline
\textbf{Tecnologia} & \textbf{Finalidade no projeto} \\ \hline
\endfirsthead
\hline
\textbf{Tecnologia} & \textbf{Finalidade no projeto} \\ \hline
\endhead
Python 3.12 (compatível com 3.10 ou superior) & Linguagem de implementação \\ \hline
Tkinter & Interface gráfica (janela, palco do avatar, painel e configurações) \\ \hline
Pillow e NumPy \cite{harris2020} & Recorte, nitidez, troca de roupa e geração dos quadros do avatar; análise de áudio \\ \hline
SpeechRecognition e PyAudio & Captura do microfone e reconhecimento de fala pelo serviço do Google \\ \hline
Vosk (modelo \cod{vosk-model-small-pt-0.3}) & Reconhecimento local da frase de ativação, com gramática restrita \\ \hline
edge-tts e gTTS & Síntese de voz neural com ajuste de tom e ritmo; voz de reserva \\ \hline
pygame & Decodificação e reprodução do áudio em memória \\ \hline
SDK oficial da Anthropic (Claude, modelo \cod{claude-opus-5-5}) & Compreensão de linguagem natural com ferramentas e transmissão em fluxo contínuo \\ \hline
SQLite & Agenda, histórico auditável e preferências \\ \hline
pytest e pytest-cov & Testes automatizados e medição de cobertura \\ \hline
ruff & Análise estática de estilo, de defeitos comuns e de complexidade \\ \hline
mypy & Verificação estática de tipos \\ \hline
pre-commit e GitHub Actions com Xvfb & Verificações antes de cada versão e integração contínua, inclusive da interface gráfica \\ \hline
dpkg-deb e fakeroot & Construção do pacote de instalação para Ubuntu \\ \hline
GNU Make & Automação das tarefas de desenvolvimento \\ \hline
\end{tabularx}
\fonte{elaboração própria (2026).}
\end{singlespace}

O ambiente de desenvolvimento e de testes foi um computador com Ubuntu 24.04 (núcleo Linux 6.8),
oito núcleos de processamento, microfone interno e um microfone USB, com Python 3.12 em ambiente
virtual isolado.

\section{Critérios de validação}

O sistema foi considerado válido quando atendeu, simultaneamente, aos critérios a seguir:

\begin{enumerate}[label=\alph*)]
  \item todos os testes automatizados aprovados, inclusive os de interface gráfica e os de
        arquitetura;
  \item nenhuma ocorrência apontada pela análise estática de estilo e de complexidade, com
        complexidade ciclomática máxima de 10 por função, nem pela verificação de tipos;
  \item cada requisito funcional associado a pelo menos um teste automatizado ou a um procedimento de
        verificação manual documentado (Apêndice A);
  \item metas de desempenho atingidas, conforme os requisitos não funcionais;
  \item instalador testado em ambiente isolado, sem alterar o sistema;
  \item execução real do assistente pelo usuário, com as funções principais operando.
\end{enumerate}

Para a revisão segundo os princípios SOLID, adotou-se o critério de preservação do comportamento: a
suíte completa deveria permanecer aprovada após cada passo de refatoração e, quando a alteração
afetasse a geração de imagens, a saída deveria permanecer idêntica pixel a pixel.

\section{Normalização do documento}

Este documento segue as normas da Associação Brasileira de Normas Técnicas para trabalhos acadêmicos
\cite{abnt14724}, referências \cite{abnt6023} e citações \cite{abnt10520}. Por se tratar de uma
documentação exclusivamente textual, os fluxos e as estruturas do sistema são descritos em prosa, em
listas numeradas e em quadros, sem ilustrações gráficas.

\chapter{Requisitos do sistema}

Este capítulo especifica os requisitos do Robofael. Seguindo \textcite{sommerville2018}, os
requisitos funcionais (RF) descrevem os serviços que o sistema oferece, e os requisitos não funcionais
(RNF) descrevem restrições e atributos de qualidade, organizados segundo as características da
ISO/IEC 25010 \cite{iso25010}. Cada requisito recebeu um identificador, uma prioridade e a indicação
de como é verificado. A prioridade segue três níveis: \textbf{essencial}, sem o qual o sistema não
cumpre seu propósito; \textbf{importante}, que afeta significativamente a experiência; e
\textbf{desejável}, que agrega valor sem comprometer o uso básico. A rastreabilidade entre os
requisitos e os testes automatizados é apresentada no Apêndice A.

\section{Partes interessadas e atores}

O sistema possui um único ator humano, o \textbf{usuário}, que é a pessoa dona do computador e das
fotografias que dão origem ao avatar. Ele interage por voz, teclado e mouse. Há ainda três atores
externos: o \textbf{serviço de modelo de linguagem} (Claude, na nuvem, ou um modelo local servido pelo
Ollama), que interpreta pedidos em linguagem natural; os \textbf{serviços de fala} (reconhecimento do Google e síntese do edge-tts e do gTTS),
acessados pela internet; e o \textbf{sistema operacional}, que fornece o microfone, o áudio, a lista de
programas instalados e os diretórios de dados do usuário.

\section{Requisitos funcionais}

Os requisitos funcionais foram agrupados por área: interação e avatar (Quadro~\ref{qua:rf-interacao}),
agenda e memória (Quadro~\ref{qua:rf-agenda}), inteligência e integração com o sistema
(Quadro~\ref{qua:rf-inteligencia}) e configuração, idioma e distribuição (Quadro~\ref{qua:rf-config}).
Os requisitos incorporados nas versões 3.6.0 a 3.8.0 estão reunidos no Quadro~\ref{qua:rf-novos}.

\begin{singlespace}
\legendaquadro{Requisitos funcionais de interação e avatar}\label{qua:rf-interacao}
\small
\begin{tabularx}{\textwidth}{|C{1.25cm}|L|C{2.0cm}|P{3.2cm}|}
\hline
\textbf{ID} & \textbf{Descrição} & \textbf{Prioridade} & \textbf{Verificação} \\ \hline
\endfirsthead
\hline
\textbf{ID} & \textbf{Descrição} & \textbf{Prioridade} & \textbf{Verificação} \\ \hline
\endhead
RF01 & Exibir um avatar em formato de busto, centralizado, gerado a partir das fotografias do usuário, com
       roupa substituída por uma camisa de estilo futurista e nitidez realçada. & Essencial & Testes de avatar \\ \hline
RF02 & Animar a boca do avatar em sincronia com a voz sintetizada, variando abertura (12 níveis) e
       formato (três formatos de vogal), com transição suave e leve balanço em repouso. & Essencial & Testes de animação \\ \hline
RF03 & Iniciar a escuta de um comando pela tecla ESPAÇO ou pelo botão de microfone, exibindo o nível
       real da voz captada. & Essencial & Testes do orquestrador; uso real \\ \hline
RF04 & Ativar o assistente sem as mãos ao ouvir a frase “ok Fael”, confirmando com um sinal sonoro e
       passando a ouvir o comando. & Essencial & Testes de ativação \\ \hline
RF05 & Aceitar a frase de ativação e o comando em uma única fala (por exemplo, “ok Fael, anota tarefa
       comprar pão”), extraindo apenas o comando. & Importante & Testes de ativação (transcrições reais) \\ \hline
RF06 & Responder por voz neural, com tom e ritmo ajustáveis, inclusive a partir do perfil gravado da
       voz do usuário. & Essencial & Testes de voz e configurações \\ \hline
RF07 & Aceitar comandos digitados em uma caixa de texto, com o mesmo tratamento dos comandos falados; um texto
       enviado com o assistente ocupado interrompe o que ele faz e é atendido em seguida. & Importante & Testes de interface e do orquestrador \\ \hline
RF08 & Exibir legendas com a fala reconhecida do usuário e com a resposta à medida que ela é gerada. & Importante & Testes de interface \\ \hline
RF09 & Permitir interromper, a qualquer momento, a escuta, a geração ou a fala da resposta (ESPAÇO, ESC
       ou botão contextual). & Essencial & Testes do orquestrador \\ \hline
RF10 & Indicar visualmente o estado do assistente (iniciando, parado, ouvindo, pensando, falando), com
       rótulo do botão adequado a cada estado. & Importante & Testes de estados e de interface \\ \hline
RF11 & Saudar o usuário ao iniciar, informando o número de tarefas pendentes, o modo offline (se for o
       caso) e como chamar o assistente. & Desejável & Testes de sessão e de idiomas \\ \hline
RF12 & Exibir avisos na tela para microfone ausente, voz captada baixa demais, falhas de áudio e de
       avatar e desligamento da escuta passiva. & Importante & Testes de ganho, ativação e interface \\ \hline
RF13 & Oferecer modo contínuo, em que o assistente volta a ouvir após cada resposta até que o usuário se
       despeça. & Desejável & Testes do orquestrador \\ \hline
RF14 & Trocar o avatar entre as fotografias disponíveis (tecla T ou configurações), sem travar a interface. & Desejável & Testes de configurações \\ \hline
\end{tabularx}
\fonte{elaboração própria (2026).}
\end{singlespace}

\begin{singlespace}
\legendaquadro{Requisitos funcionais de agenda e memória}\label{qua:rf-agenda}
\small
\begin{tabularx}{\textwidth}{|C{1.25cm}|L|C{2.0cm}|P{3.2cm}|}
\hline
\textbf{ID} & \textbf{Descrição} & \textbf{Prioridade} & \textbf{Verificação} \\ \hline
\endfirsthead
\hline
\textbf{ID} & \textbf{Descrição} & \textbf{Prioridade} & \textbf{Verificação} \\ \hline
\endhead
RF15 & Criar tarefas por voz ou por texto. & Essencial & Testes da agenda e do modo offline \\ \hline
RF16 & Listar as tarefas pendentes e, opcionalmente, as concluídas. & Essencial & Testes da agenda e de ferramentas \\ \hline
RF17 & Concluir, reabrir e remover tarefas por voz ou pelo painel (duplo clique e tecla Delete). & Essencial & Testes da agenda e de interface \\ \hline
RF18 & Criar lembretes para uma data e hora ou para daqui a N minutos e anunciá-los por voz no horário. & Essencial & Testes do agendador \\ \hline
RF19 & Listar e remover lembretes futuros. & Importante & Testes da agenda \\ \hline
RF20 & Registrar automaticamente, em histórico, cada criação, conclusão, reabertura e remoção de tarefa e
       cada criação, aviso e remoção de lembrete, com data e hora. & Essencial & Testes do SQLite \\ \hline
RF21 & Responder “o que eu fiz” em um dia, semana ou período (ontem, anteontem, um dia da semana, uma data,
       semana passada, esta semana, últimos N dias, mês passado), considerando apenas o resultado
       líquido (tarefa reaberta não conta como feita). & Essencial & Testes de memória \\ \hline
RF22 & Verificar a integridade do histórico e apontar o primeiro registro alterado, removido ou inserido
       indevidamente (comando \cod{make integridade}). & Importante & Testes do SQLite \\ \hline
RF23 & Migrar automaticamente a agenda do formato JSON das versões anteriores para o banco de dados,
       preservando o arquivo original renomeado. & Importante & Testes do SQLite \\ \hline
\end{tabularx}
\fonte{elaboração própria (2026).}
\end{singlespace}

\begin{singlespace}
\legendaquadro{Requisitos funcionais de inteligência e integração com o sistema}\label{qua:rf-inteligencia}
\small
\begin{tabularx}{\textwidth}{|C{1.25cm}|L|C{2.0cm}|P{3.2cm}|}
\hline
\textbf{ID} & \textbf{Descrição} & \textbf{Prioridade} & \textbf{Verificação} \\ \hline
\endfirsthead
\hline
\textbf{ID} & \textbf{Descrição} & \textbf{Prioridade} & \textbf{Verificação} \\ \hline
\endhead
RF24 & Interpretar pedidos em linguagem natural com um modelo de linguagem que opera a agenda, o histórico
       e os programas por meio de ferramentas. & Essencial & Testes do modelo Claude e de ferramentas \\ \hline
RF25 & Funcionar em modo offline quando o modelo de linguagem estiver indisponível, atendendo a comandos de
       tarefas, lembretes relativos, hora, data, histórico e programas. & Essencial & Testes do modo offline \\ \hline
RF26 & Iniciar uma nova conversa quando solicitado (“nova conversa”), sem perder a agenda. & Desejável & Testes do modelo Claude \\ \hline
RF27 & Informar a hora e a data atuais. & Importante & Testes do modo offline \\ \hline
RF28 & Abrir programas instalados no Linux pelo nome ou por sinônimos (por exemplo, “navegador”,
       “calculadora”), usando o navegador padrão quando pedido e solicitando escolha em caso de
       ambiguidade. & Importante & Testes de programas \\ \hline
RF29 & Listar os programas instalados, com filtro opcional pelo nome. & Desejável & Testes de programas \\ \hline
\end{tabularx}
\fonte{elaboração própria (2026).}
\end{singlespace}

\begin{singlespace}
\legendaquadro{Requisitos funcionais de configuração, idioma e distribuição}\label{qua:rf-config}
\small
\begin{tabularx}{\textwidth}{|C{1.25cm}|L|C{2.0cm}|P{3.2cm}|}
\hline
\textbf{ID} & \textbf{Descrição} & \textbf{Prioridade} & \textbf{Verificação} \\ \hline
\endfirsthead
\hline
\textbf{ID} & \textbf{Descrição} & \textbf{Prioridade} & \textbf{Verificação} \\ \hline
\endhead
RF30 & Gravar amostras da voz do usuário e gerar um perfil com o deslocamento de tom e de ritmo que
       aproxima a voz sintética da voz dele. & Importante & Testes de perfil de voz \\ \hline
RF31 & Oferecer tela de configurações com idioma, avatar, voz base, tom, ritmo, botão para ouvir o
       resultado, uso do perfil gravado e ativação por voz, com opções de salvar e de cancelar (que
       desfaz o teste). & Importante & Testes de configurações \\ \hline
RF32 & Guardar as preferências entre sessões, com precedência sobre o perfil gravado e sobre os valores
       padrão. & Importante & Testes de configurações \\ \hline
RF33 & Responder em português do Brasil ou em inglês dos Estados Unidos, trocando ao vivo a interface,
       as respostas, a voz, o reconhecimento de fala e as instruções do modelo. & Importante & Testes de idiomas \\ \hline
RF34 & Entender comandos nos dois idiomas, independentemente do idioma das respostas. & Desejável & Testes de idiomas \\ \hline
RF35 & Aceitar opções de linha de comando (avatar, modo contínuo, desligar ativação, idioma e registros
       detalhados) e configuração por variáveis de ambiente. & Desejável & Testes de configuração \\ \hline
RF36 & Ser instalável no Ubuntu por um pacote \cod{.deb}, com comandos no terminal, atalhos e ícone no
       menu de aplicativos e remoção que preserva os dados do usuário. & Importante & Teste do pacote em ambiente isolado \\ \hline
\end{tabularx}
\fonte{elaboração própria (2026).}
\end{singlespace}

\begin{singlespace}
\legendaquadro{Requisitos funcionais incorporados nas versões 3.6.0 a 3.8.0}\label{qua:rf-novos}
\small
\begin{tabularx}{\textwidth}{|C{1.25cm}|L|C{2.0cm}|P{3.2cm}|}
\hline
\textbf{ID} & \textbf{Descrição} & \textbf{Prioridade} & \textbf{Verificação} \\ \hline
\endfirsthead
\hline
\textbf{ID} & \textbf{Descrição} & \textbf{Prioridade} & \textbf{Verificação} \\ \hline
\endhead
RF37 & Sem o Claude, conversar e atender pedidos de agenda em linguagem livre com um modelo local
       (Ollama), usando apenas ferramentas que criam ou consultam, e voltar ao modo offline se o modelo
       local deixar de responder. & Importante & Testes do modelo local \\ \hline
RF38 & Exibir nas configurações o modelo de linguagem em uso e os modelos compatíveis do Claude e do
       Ollama, informando o motivo quando um provedor não puder ser consultado. & Desejável & Testes de banco e modelos; testes de configurações \\ \hline
RF39 & Escolher o microfone entre os conectados ao computador, com a troca valendo na escuta seguinte,
       inclusive para a escuta passiva e a gravação da voz. & Importante & Testes de microfone e inicialização \\ \hline
RF40 & Oferecer, como opções, abrir o assistente ao entrar no computador e pelo comando
       \cod{robofael} em qualquer terminal. & Desejável & Testes de microfone e inicialização \\ \hline
RF41 & Mudar o banco de dados para outra pasta com o assistente em uso, sem apagar o arquivo anterior e
       sem sobrescrever um banco existente, mantendo a escolha nas aberturas seguintes. & Importante & Testes de banco e modelos \\ \hline
RF42 & Fazer backup do banco de dados em qualquer local, com a cópia conferida antes de ser aceita. & Importante & Testes de banco e modelos \\ \hline
\end{tabularx}
\fonte{elaboração própria (2026).}
\end{singlespace}

\section{Requisitos não funcionais}

Os requisitos não funcionais estão organizados em duas partes: desempenho, confiabilidade e segurança
(Quadro~\ref{qua:rnf-1}); e usabilidade, manutenibilidade e portabilidade (Quadro~\ref{qua:rnf-2}).
Sempre que possível, foram definidos com metas mensuráveis.

\begin{singlespace}
\legendaquadro{Requisitos não funcionais de desempenho, confiabilidade e segurança}\label{qua:rnf-1}
\small
\begin{tabularx}{\textwidth}{|C{1.4cm}|P{3.1cm}|L|P{2.8cm}|}
\hline
\textbf{ID} & \textbf{Característica} & \textbf{Descrição e meta} & \textbf{Verificação} \\ \hline
\endfirsthead
\hline
\textbf{ID} & \textbf{Característica} & \textbf{Descrição e meta} & \textbf{Verificação} \\ \hline
\endhead
RNF01 & Desempenho (tempo) & A janela deve ficar visível em menos de 1 s após a abertura; o trabalho lento
        (modelo, microfone, áudio e avatar) ocorre em segundo plano. & Medição; teste de inicialização \\ \hline
RNF02 & Desempenho (tempo) & A primeira palavra da resposta deve ser falada assim que a primeira frase
        for gerada e sintetizada, sem esperar o restante da resposta. & Teste que prova o pipeline \\ \hline
RNF03 & Desempenho (tempo) & Uma interrupção deve surtir efeito em cerca de 10 ms em qualquer estágio
        (ouvindo, pensando ou falando). & Testes do orquestrador \\ \hline
RNF04 & Desempenho (tempo) & Lembretes devem ser anunciados com atraso inferior a 0,1 s em relação ao
        horário marcado. & Testes do agendador \\ \hline
RNF05 & Desempenho (recursos) & A animação deve rodar a 60 quadros por segundo com o assistente ativo e a
        30 em repouso, com compensação de deriva, reenviando ao Tk apenas o que mudou e trocando apenas
        imagens pré-renderizadas; falas repetidas e quadros do avatar devem vir de
        cache em disco (LRU de 300 falas; cache do avatar invalidado pelo resumo do código). & Testes de animação e de voz \\ \hline
RNF06 & Confiabilidade (integridade) & Toda alteração da agenda e seus registros de histórico devem ser
        gravados na mesma transação; o histórico é somente de acréscimo e encadeado por SHA-256. & Testes do SQLite \\ \hline
RNF07 & Confiabilidade (consistência) & O histórico da conversa com o modelo só cresce por acréscimo;
        uma rodada que falhe, seja recusada ou interrompida é descartada por inteiro. & Testes do modelo Claude \\ \hline
RNF08 & Confiabilidade (tolerância a falhas) & Na ausência do modelo, usar o modo offline; na falha da voz
        neural, usar a voz de reserva; sem microfone, aceitar texto; após três falhas seguidas da escuta
        passiva, desligá-la com aviso. & Testes de composição, ativação e voz \\ \hline
RNF09 & Confiabilidade (concorrência) & Apenas uma interação por vez; somente a thread principal acessa a
        interface gráfica; cada interação possui um token de cancelamento próprio, criado de forma síncrona. & Testes do orquestrador e do barramento \\ \hline
RNF10 & Segurança (privacidade) & A escuta passiva deve ser processada localmente; o áudio só é enviado a
        serviços externos depois da ativação ou de um comando explícito. Os dados ficam em arquivos locais
        do usuário. & Revisão de código; testes de ativação \\ \hline
RNF11 & Segurança (integridade de entradas) & As entradas das ferramentas do modelo devem ser validadas
        contra o esquema (tipos, campos obrigatórios e campos extras) antes da execução; entradas truncadas
        nunca são executadas; programas são abertos sem interpretador de comandos. & Testes de ferramentas e de programas \\ \hline
RNF12 & Segurança (segredos) & Credenciais vêm de variáveis de ambiente e não são registradas em logs nem
        gravadas em arquivos do projeto. & Revisão de código \\ \hline
RNF13 & Exatidão (ativação) & A frase de ativação não deve disparar com frases parecidas: meta de zero
        falsos positivos no conjunto de calibração, inclusive para o nome do usuário (“ok Rafael”). & Testes de ativação com transcrições reais \\ \hline
\end{tabularx}
\fonte{elaboração própria (2026).}
\end{singlespace}

\begin{singlespace}
\legendaquadro{Requisitos não funcionais de usabilidade, manutenibilidade e portabilidade}\label{qua:rnf-2}
\small
\begin{tabularx}{\textwidth}{|C{1.4cm}|P{3.1cm}|L|P{2.8cm}|}
\hline
\textbf{ID} & \textbf{Característica} & \textbf{Descrição e meta} & \textbf{Verificação} \\ \hline
\endfirsthead
\hline
\textbf{ID} & \textbf{Característica} & \textbf{Descrição e meta} & \textbf{Verificação} \\ \hline
\endhead
RNF14 & Usabilidade (operabilidade) & Todas as ações principais devem ter atalho de teclado (ESPAÇO, ENTER,
        ESC, T, Ctrl+Q); o botão principal muda de rótulo conforme o estado. & Testes de interface \\ \hline
RNF15 & Usabilidade (estética e estabilidade) & A janela tem tamanho fixo, para que o palco desenhado em
        coordenadas absolutas não se desloque ao maximizar. & Teste de tamanho fixo \\ \hline
RNF16 & Usabilidade (retorno) & O nível do microfone, o estado atual, as legendas e os avisos devem ser
        exibidos continuamente. & Testes de interface \\ \hline
RNF17 & Usabilidade (inclusão) & As instruções ao modelo e as mensagens não devem atribuir gênero ao
        usuário. & Revisão do catálogo de mensagens \\ \hline
RNF18 & Internacionalização & Todas as mensagens ficam em um catálogo único, com as duas línguas lado a lado;
        nenhuma tradução pode faltar e os campos de formatação devem coincidir. & Testes de idiomas \\ \hline
RNF19 & Manutenibilidade (modularidade) & Arquitetura hexagonal com regra de dependência verificada por
        teste; domínio e portas sem bibliotecas de infraestrutura. & Testes de arquitetura \\ \hline
RNF20 & Manutenibilidade (analisabilidade) & Complexidade ciclomática de no máximo 10 por função;
        funções curtas; nomes padronizados; nenhum parâmetro sem uso; tipos verificados estaticamente. & ruff e mypy na integração contínua \\ \hline
RNF21 & Manutenibilidade (testabilidade) & Cobertura mínima de 80\% das linhas e ramificações; dublês para
        todos os serviços externos; testes de interface com servidor gráfico virtual. & pytest-cov; integração contínua \\ \hline
RNF22 & Manutenibilidade (modificabilidade) & Acrescentar uma intenção offline, uma ferramenta do modelo ou
        uma expressão de período não deve exigir alteração do código de despacho. & Testes de arquitetura \\ \hline
RNF23 & Portabilidade (adaptabilidade) & Compatível com Python 3.10 ou superior; caminhos de dados
        configuráveis e, quando instalado, nos diretórios XDG do usuário. & Integração contínua (3.10 e 3.12) \\ \hline
RNF24 & Portabilidade (instalabilidade) & Instalação sem internet, com as dependências Python incluídas no
        pacote; o ambiente Python é criado na instalação e verificado com \cod{pip check}. & Teste do pacote \\ \hline
RNF25 & Rastreabilidade de mudanças & Versões numeradas segundo o versionamento semântico, com registro de
        mudanças por versão. & Revisão do CHANGELOG \\ \hline
RNF26 & Desempenho (recursos) & Em repouso, o assistente deve consumir no máximo 10\% de um núcleo de
        processamento; um quadro em repouso não pode reenviar ao Tk o que não mudou. & Medição;
        \cod{test_renderizacao} \\ \hline
RNF27 & Desempenho (fluidez) & Nenhum trabalho da thread da interface pode ultrapassar o orçamento de um
        quadro; conversões de imagem são feitas aos poucos. Entre frases faladas, a pausa não deve exceder
        0,2 s. & \cod{test_renderizacao}; testes de áudio \\ \hline
RNF28 & Confiabilidade (preservação dos dados) & Nenhuma operação sobre o banco pode apagar ou sobrescrever
        dados: cópias são conferidas antes do uso, o arquivo anterior é renomeado, e um banco escolhido
        inacessível não é substituído por um banco vazio. Os testes não podem alterar a pasta pessoal
        real. & Testes de banco e modelos; fixture de pasta pessoal isolada \\ \hline
\end{tabularx}
\fonte{elaboração própria (2026).}
\end{singlespace}

\section{Regras de negócio}

Além dos requisitos, algumas regras de negócio orientam o comportamento do sistema:

\begin{enumerate}[label=RN\arabic*:,leftmargin=2.2cm]
  \item comandos de agenda têm prioridade sobre a abertura de programas: “anota tarefa abrir conta no
        banco” cria uma tarefa, e não tenta abrir um programa chamado “conta”;
  \item períodos ditos pelo usuário referem-se sempre ao passado: “segunda-feira” é a mais recente antes
        de hoje, e “dia 30”, dito em um dia 22, é o dia 30 do mês anterior;
  \item datas numéricas seguem o idioma das respostas: 05/10 é 5 de outubro em português e 10 de maio em
        inglês;
  \item para responder “o que eu fiz”, vale o estado líquido no período: uma tarefa concluída e depois
        reaberta não é contada como feita;
  \item a precedência das configurações é: preferência salva, depois opção de linha de comando ou
        variável de ambiente e, por fim, valor padrão; no caso da voz, o perfil gravado vem antes do
        padrão;
  \item ao trocar de idioma, a voz é substituída pela voz padrão do novo idioma caso a atual não o fale,
        e a conversa com o modelo recomeça;
  \item nenhum dado do usuário é apagado automaticamente: arquivos corrompidos, migrados ou substituídos
        por uma mudança de local do banco são renomeados e preservados, e a remoção do pacote mantém a
        pasta de dados;
  \item com o modelo local, ações que alteram ou removem dados e a abertura de programas só acontecem
        por comandos explícitos, nunca por escolha do modelo.
\end{enumerate}

\chapter{Arquitetura e projeto}

Este capítulo descreve a organização do Robofael: as camadas e seus módulos, os contratos entre elas,
o modelo de dados, o modelo de concorrência e os principais fluxos de execução. Por se tratar de uma
documentação textual, os fluxos são apresentados como sequências numeradas de passos.

\section{Visão geral}

O Robofael é uma aplicação de área de trabalho escrita em Python, distribuída como o pacote
\cod{robofael} e executada como um único processo. A arquitetura segue o
modelo hexagonal \cite{cockburn2005}: no centro estão o domínio e os casos de uso; nas bordas, os
adaptadores para cada tecnologia; e um único ponto, a raiz de composição, conhece as implementações
concretas e monta o sistema. A regra de dependência \cite{martin2019} é respeitada em todo o código e
verificada por testes automatizados (Seção~\ref{sec:testes-arquitetura}).

\section{Camadas e módulos}

O Quadro~\ref{qua:camadas} descreve as camadas, os módulos que as compõem e as dependências
permitidas a cada uma.

\begin{singlespace}
\legendaquadro{Camadas da arquitetura, módulos e dependências permitidas}\label{qua:camadas}
\small
\begin{tabularx}{\textwidth}{|P{3.0cm}|L|P{3.4cm}|}
\hline
\textbf{Camada} & \textbf{Módulos e responsabilidade} & \textbf{Pode depender de} \\ \hline
\endfirsthead
\hline
\textbf{Camada} & \textbf{Módulos e responsabilidade} & \textbf{Pode depender de} \\ \hline
\endhead
Domínio (\cod{dominio}) & Entidades e regras puras: \cod{modelos} (tarefa, lembrete, registro de
auditoria), \cod{agenda} (serviço de domínio transacional), \cod{repositorios} (contrato de
persistência exigido pela agenda), \cod{estados} (máquina de estados), \cod{eventos}, \cod{linguagem}
(despedidas, pedidos de nova conversa, perguntas sobre o histórico), \cod{periodos} (interpretação de
períodos), \cod{programas} (busca do programa pedido) e \cod{ativacao} (gramática e busca fonética da
frase de ativação). & Somente \cod{i18n} \\ \hline
Portas (\cod{portas}) & Contratos (\cod{Protocol}) pequenos: modelo de linguagem, reconhecedor de
fala, captura de áudio, detector de ativação, sintetizador (simples e configurável), reprodutor de
áudio, consulta ao histórico, preferências, arquivo do banco (mudar de lugar e copiar), persistência
completa, lançador de programas, dispositivos de áudio, inicialização com o sistema e catálogo de modelos
de linguagem. & Domínio \\ \hline
Aplicação (\cod{aplicacao}) & Casos de uso e políticas: \cod{orquestrador} (ciclo de conversa),
\cod{pipeline} (texto em fala, em fluxo contínuo), \cod{vigia} (escuta passiva), \cod{sessao}
(políticas de voz, idioma, ativação, microfone e saudação; mudança e backup do banco; modelos
disponíveis), \cod{memoria} (“o que eu fiz”), \cod{programas},
\cod{agendador} (lembretes), \cod{barramento} (eventos entre threads) e \cod{frases} (divisor
incremental de frases). & Domínio, portas, \cod{i18n} e \cod{configuracao} \\ \hline
Adaptadores (\cod{adaptadores}) & Tecnologias concretas: \cod{claude} (modelo com ferramentas e
catálogo de modelos), \cod{ollama} (modelo local com respostas em camadas e catálogo dos modelos
instalados),
\cod{ferramentas} (esquemas, validação e execução), \cod{offline} (intenções locais), \cod{sqlite}
(memória de longo prazo), \cod{microfone} (captura e reconhecimento), \cod{ativacao_vosk},
\cod{sintetizador} (edge-tts e gTTS com cache), \cod{reprodutor} (pygame), \cod{programas_linux}
(entradas \cod{.desktop}), \cod{dispositivos_audio} (microfones do servidor de áudio),
\cod{inicializacao_linux} (abertura com o computador e comando no terminal), \cod{alsa}
(silenciamento de mensagens nativas) e \cod{repositorio_json}
(leitura do formato legado). & Domínio, portas, aplicação, \cod{audio} e \cod{i18n} \\ \hline
Núcleos de processamento (\cod{audio}, \cod{avatar}) & Processamento de sinais (tom, envelope de
volume e timbre, detecção de fala, ganho automático, perfil de voz) e geração das imagens do avatar
(catálogo, renderização paralela com cache e suavização da animação). & \cod{configuracao} e \cod{i18n} \\ \hline
Apresentação (\cod{apresentacao}) & Interface Tkinter: \cod{app} (janela principal), \cod{palco}
(avatar, anéis de estado e legendas), \cod{cena} (canvas diferencial), \cod{laco} (laço de quadros
adaptativo e fila de trabalho com orçamento), \cod{painel} (agenda), \cod{configuracoes} (diálogo),
\cod{gravador} (ferramenta de gravação da voz) e \cod{tema}. & Aplicação, domínio, portas, núcleos e
raiz de composição \\ \hline
Raiz de composição (\cod{bootstrap}) & Escolhe e conecta as implementações concretas; inicia os
serviços lentos em paralelo. & Todas as camadas internas e os adaptadores \\ \hline
Transversais (\cod{i18n}, \cod{configuracao}) & Catálogo de mensagens e tradutor; configuração
imutável lida do ambiente. & Nenhuma camada de infraestrutura \\ \hline
\end{tabularx}
\fonte{elaboração própria (2026).}
\end{singlespace}

\section{Contratos entre as camadas}

Os contratos são definidos como protocolos estruturais \cite{pep544}. Um componente satisfaz um
contrato quando possui os métodos e atributos exigidos, sem precisar herdar de uma classe base. Isso
permite que os testes usem dublês simples \cite{meszaros2007} e que a substituição de uma tecnologia,
como trocar o serviço de síntese de voz, não afete o domínio nem a aplicação. Os contratos foram
mantidos pequenos, conforme o princípio da segregação de interfaces. Por exemplo, o serviço de memória
depende apenas da consulta ao histórico; a agenda, apenas do repositório da agenda; e a tela de
configurações, apenas de um sintetizador cuja voz possa ser trocada.

O contrato de persistência da agenda pertence ao próprio domínio (\cod{dominio/repositorios.py}),
pois é o domínio que declara do que precisa. A infraestrutura, representada pela classe
\cod{BancoSqlite}, implementa esse contrato e também os de histórico, de preferências e de arquivo
do banco. A porta \cod{PersistenciaAssistente} combina os quatro contratos para os poucos pontos que
precisam do conjunto, como o núcleo da sessão.

Os modelos de linguagem informam o provedor e o modelo em uso, e cada provedor oferece um catálogo
(\cod{CatalogoModelos}) com os modelos que o assistente consegue usar. Quando um catálogo não pode ser
consultado, ele levanta uma exceção com a chave da mensagem no catálogo de textos, e não um texto
pronto: o adaptador não precisa conhecer o idioma da interface.

\section{Modelo de dados}

A persistência utiliza um único arquivo SQLite (\cod{memoria.db}), cujo esquema é versionado por meio
do mecanismo \cod{PRAGMA user_version} e migrado automaticamente na abertura \cite{sqlite}. As tabelas
são descritas no Quadro~\ref{qua:tabelas}.

\begin{singlespace}
\legendaquadro{Tabelas do banco de dados}\label{qua:tabelas}
\small
\begin{tabularx}{\textwidth}{|P{2.9cm}|L|}
\hline
\textbf{Tabela} & \textbf{Conteúdo e restrições} \\ \hline
\endfirsthead
\hline
\textbf{Tabela} & \textbf{Conteúdo e restrições} \\ \hline
\endhead
\cod{tarefas} & Identificador, texto, data de criação, indicador de conclusão (restrito a 0 ou 1) e
data de conclusão. \\ \hline
\cod{lembretes} & Identificador, texto, data e hora do aviso e indicador de aviso já dado. \\ \hline
\cod{historico} & Sequência, data e hora do fato, tipo do evento, entidade, identificador da entidade,
descrição no momento do fato, dados adicionais em JSON, resumo do registro anterior e resumo do
próprio registro (único). Índice por data. Gatilhos recusam qualquer UPDATE ou DELETE. \\ \hline
\cod{preferencias} & Pares chave e valor com as escolhas da tela de configurações (avatar, voz, tom,
ritmo, ativação por voz, idioma e microfone). \\ \hline
\cod{meta} & Pares chave e valor de controle interno, como o próximo identificador. \\ \hline
\end{tabularx}
\fonte{elaboração própria (2026).}
\end{singlespace}

Os tipos de evento registrados no histórico são: tarefa criada, tarefa concluída, tarefa reaberta,
tarefa removida, lembrete criado, lembrete avisado e lembrete removido. Cada registro guarda a
descrição do item no momento do fato, de modo que a resposta sobre o passado continua correta mesmo
se o item for removido depois. O resumo de cada registro é calculado com SHA-256 \cite{nist2015} sobre
o resumo anterior e os campos do registro; a verificação de integridade recalcula a cadeia e aponta o
primeiro registro inconsistente, detectando alteração, remoção, inserção por fora do sistema e
quebra de encadeamento.

O arquivo fica, por padrão, na pasta de dados, mas pode ser movido para outra pasta pela tela de
configurações. Como o banco não pode guardar o próprio endereço, o local escolhido é registrado em um
pequeno arquivo JSON na pasta de dados, lido na abertura. A mudança e o backup usam a interface de
backup do SQLite \cite{sqlite}, que produz uma cópia consistente mesmo com gravações em andamento; a
cópia é conferida com \cod{PRAGMA integrity_check} antes de ser aceita, e o arquivo anterior é
renomeado com a data, nunca apagado.

\section{Modelo de concorrência}

Um assistente de voz precisa ouvir, pensar, falar e animar ao mesmo tempo, mas a biblioteca Tkinter
não é segura para uso por várias threads. A solução adotada combina quatro mecanismos:

\begin{enumerate}[label=\alph*)]
  \item \textbf{barramento de eventos}: as threads de trabalho publicam eventos imutáveis (mudança de
        estado, fala reconhecida, resposta parcial, nível do microfone, aviso, serviços prontos, avatar
        pronto) em uma fila segura; a thread principal os despacha a cada quadro da animação, e somente
        ela altera componentes gráficos;
  \item \textbf{uma interação por vez}: uma trava garante que apenas uma interação (ouvir, pensar e
        falar) esteja em andamento; anúncios, como saudação e lembretes, entram em uma fila e são falados
        entre as interações;
  \item \textbf{token de cancelamento por interação}: um objeto \cod{threading.Event} novo é criado de
        forma síncrona no início de cada interação; todos os estágios consultam esse token, de modo que
        uma interrupção nunca se perde e estágios cancelados anteriormente não são reativados;
  \item \textbf{máquina de estados explícita}: as transições válidas são declaradas (por exemplo, de
        “ouvindo” só é possível ir para “parado” ou “pensando”); uma transição inválida gera erro
        imediato, em vez de um defeito silencioso.
\end{enumerate}

As threads existentes em execução normal são: a principal (interface e animação); a de inicialização
dos serviços; as de renderização do avatar e de geração dos halos; a da interação corrente; as dos
estágios “pensar” e “sintetizar” do pipeline, com duas threads de síntese de voz; a da escuta passiva;
e a do agendador de lembretes, que dorme até o próximo vencimento.

Como só a thread principal pode criar imagens para a biblioteca gráfica, a conversão das imagens
preparadas em segundo plano (quadros do avatar e halos) passa por uma fila de trabalho com orçamento de
4 ms por quadro. Cada tarefa é um gerador que converte uma imagem e devolve o controle, de modo que um
trabalho de centenas de milissegundos se espalha por vários quadros sem interromper a animação.

\section{Fluxos de execução}

\subsection{Inicialização}

\begin{enumerate}[label=\arabic*.]
  \item O ponto de entrada lê as opções da linha de comando e a configuração imutável do ambiente.
  \item A raiz de composição monta o núcleo, uma operação instantânea: abre o banco de dados no local
        escolhido ou no padrão (migrando a agenda JSON antiga, se houver), cria a agenda, carrega as preferências, define o idioma e cria
        os serviços de memória e de programas.
  \item A janela é exibida imediatamente, com tamanho fixo, e o laço de quadros começa; ele roda a 60
        quadros por segundo com o assistente ativo e a 30 em repouso.
  \item Em segundo plano, o avatar é renderizado ou lido do cache, os halos são gerados, e os serviços
        lentos são iniciados
        em paralelo: conexão com o modelo de linguagem, tentando o Claude, depois o modelo local
        (Ollama) e, por fim, o modo offline;
        calibração do microfone; abertura do áudio; carregamento do modelo de ativação; e leitura da
        lista de programas instalados.
  \item Quando os serviços ficam prontos, um evento avisa a interface, que habilita o botão de
        microfone, inicia o agendador, enfileira a saudação e liga a escuta passiva.
\end{enumerate}

\subsection{Conversa em pipeline}

\begin{enumerate}[label=\arabic*.]
  \item O usuário pressiona ESPAÇO, diz “ok Fael” ou envia um texto; o orquestrador adquire a trava de
        interação, libera o microfone da escuta passiva e cria um token de cancelamento.
  \item No estado “ouvindo”, o reconhecedor capta a fala com detecção de atividade de voz e devolve o
        texto; o nível do microfone é publicado continuamente para a interface.
  \item No estado “pensando”, o modelo começa a transmitir a resposta em pedaços de texto.
  \item O estágio “pensar” do pipeline acumula os pedaços, publica a resposta parcial para a legenda e
        entrega frases completas a uma fila; o divisor de frases corta a primeira frase cedo e não quebra
        números nem abreviações.
  \item O estágio “sintetizar” pede a síntese de cada frase assim que ela fica pronta, com até duas
        sínteses simultâneas; a thread que toca decodifica o áudio, apara o silêncio das pontas (a voz
        sintetizada traz cerca de 1 s de silêncio ao fim de cada frase) e calcula o envelope de volume e
        de timbre que comandará a boca do avatar.
  \item A thread da interação toca as faixas em ordem, no estado “falando”, enquanto as frases seguintes
        ainda estão sendo geradas e sintetizadas.
  \item Ao fim, a interação publica a resposta concluída, volta ao estado “parado”, fala os anúncios
        pendentes e libera a trava; a escuta passiva retoma o microfone.
\end{enumerate}

Esse arranjo faz com que a espera percebida seja a da primeira frase, e não a da resposta inteira. Uma
interrupção sinaliza o token: a reprodução para, as filas são esvaziadas sem sintetizar e o fluxo de
texto do modelo é fechado, o que descarta a rodada no histórico da conversa.

\subsection{Ativação por voz}

\begin{enumerate}[label=\arabic*.]
  \item Enquanto o assistente está parado, a vigia reserva o microfone e entrega o áudio ao Vosk, que só
        pode escolher entre as variantes aceitas de “ok Fael”, frases distratoras parecidas (como “o que
        foi” e “ok Rafael”) e um marcador de fala desconhecida.
  \item O áudio passa por um controle automático de ganho, que amplifica a voz até um nível adequado,
        com teto de 8 vezes e resistência a picos isolados; se a voz chega fraca demais de forma
        persistente, um aviso é exibido.
  \item Ao reconhecer uma variante aceita, a vigia continua gravando enquanto houver voz, para capturar
        um comando dito na mesma frase.
  \item Se houver comando, ele é transcrito com precisão pelo reconhecedor remoto e alinhado com a
        transcrição local, para remover a frase de ativação mesmo quando ela for mal transcrita; se não
        houver, um sinal sonoro de dois tons confirma a ativação e o assistente passa a ouvir.
  \item Em ambos os casos, a interação segue o fluxo de conversa descrito anteriormente.
\end{enumerate}

\subsection{Memória de longo prazo}

\begin{enumerate}[label=\arabic*.]
  \item Cada alteração na agenda gera, na mesma transação, um ou mais registros de auditoria.
  \item Quando o usuário pergunta “o que eu fiz…”, o período é interpretado por uma lista ordenada de
        interpretadores, do mais específico para o mais genérico, sempre voltada ao passado.
  \item O serviço de memória consulta os registros do período e calcula o resultado líquido de cada
        tarefa: concluídas, criadas e ainda pendentes, e lembretes avisados.
  \item A resposta é falada no idioma configurado; com o modelo de linguagem ativo, a mesma consulta
        está disponível como ferramenta.
\end{enumerate}

\subsection{Abertura de programas}

\begin{enumerate}[label=\arabic*.]
  \item O lançador lê as entradas \cod{.desktop} dos diretórios padronizados pelo freedesktop.org
        \cite{freedesktop}, inclusive de Snap e Flatpak, ignorando entradas ocultas e as que exigem
        programas ausentes.
  \item O pedido é extraído da fala (“abre o navegador”), normalizado e comparado com nomes, palavras-chave
        e sinônimos em português e inglês; “navegador” resolve para o navegador padrão do sistema.
  \item Se houver um candidato claro, o programa é aberto sem interpretador de comandos; se houver
        empate, o assistente pergunta qual das opções o usuário deseja.
\end{enumerate}

\subsection{Laço de ferramentas do modelo de linguagem}

\begin{enumerate}[label=\arabic*.]
  \item A fala é acrescentada ao histórico da conversa com a data e a hora atuais, mantendo as instruções
        de sistema estáveis e, portanto, reaproveitáveis pelo cache do serviço.
  \item O texto da resposta é repassado ao pipeline conforme chega.
  \item Ao fim de cada rodada, se o modelo pediu ferramentas, cada entrada é validada contra o esquema e
        executada; os resultados são devolvidos ao modelo, até o limite de oito rodadas.
  \item Recusas, entradas truncadas pelo limite de tamanho, erros da API e interrupções descartam a
        rodada inteira do histórico e produzem uma frase explicativa para o usuário.
\end{enumerate}

\subsection{Respostas do modelo local}

\begin{enumerate}[label=\arabic*.]
  \item Se a fala é um comando que o modo offline reconhece, ele responde, sem consultar o modelo.
  \item Se a fala trata de agenda, lembretes ou histórico, o assistente diz “só um instante” e envia o
        pedido ao modelo apenas com as ferramentas do assunto que criam ou consultam, nunca as que
        concluem, removem ou abrem programas.
  \item Caso contrário, a fala é conversa e vai ao modelo sem ferramentas, com o texto transmitido em
        fluxo contínuo para o pipeline.
  \item Uma chamada de ferramenta escrita como texto pelo modelo é executada, se for válida, ou
        descartada; ela nunca é falada.
  \item Se o servidor local não responde, a fala é atendida pelo modo offline, e o modelo só é tentado
        de novo depois de 30 s.
\end{enumerate}

\subsection{Animação da boca}

\begin{enumerate}[label=\arabic*.]
  \item Antes de tocar, cada faixa tem seu envelope calculado quadro a quadro: o volume define a
        abertura, e a razão entre a energia das frequências agudas (1 a 3 kHz) e graves (250 Hz a
        1 kHz) define o formato da boca (arredondada para “o” e “u”, neutra para “a” e aberta para “e” e “i”).
  \item Durante a reprodução, o reprodutor expõe a amplitude e o formato do instante atual.
  \item A cada quadro, o palco suaviza esses valores de forma independente da taxa de quadros e escolhe a
        imagem pré-renderizada correspondente, entre 36 combinações (três formatos e doze níveis).
  \item O palco descreve a cena inteira a cada quadro, mas um canvas diferencial só repassa à biblioteca
        gráfica as coordenadas e opções que mudaram, já que cada alteração obriga a redesenhar a área
        afetada, inclusive as imagens grandes que estão por baixo.
\end{enumerate}

\section{Decisões de projeto}

O Quadro~\ref{qua:decisoes} resume as principais decisões de projeto e suas justificativas.

\begin{singlespace}
\legendaquadro{Principais decisões de projeto}\label{qua:decisoes}
\small
\begin{tabularx}{\textwidth}{|C{0.8cm}|P{5cm}|L|}
\hline
\textbf{Nº} & \textbf{Decisão} & \textbf{Justificativa} \\ \hline
\endfirsthead
\hline
\textbf{Nº} & \textbf{Decisão} & \textbf{Justificativa} \\ \hline
\endhead
1 & Arquitetura hexagonal com protocolos & Troca de tecnologia sem tocar no domínio; dublês triviais nos testes. \\ \hline
2 & Pipeline de três estágios com filas & A latência percebida passa a ser a da primeira frase. \\ \hline
3 & Um token de cancelamento por interação, criado de forma síncrona & Interrupção instantânea, que nunca se perde nem reativa estágios antigos. \\ \hline
4 & Barramento de eventos despachado no laço de quadros & A interface só é alterada pela thread principal. \\ \hline
5 & Máquina de estados explícita & Estados impossíveis tornam-se erros visíveis. \\ \hline
6 & Histórico da conversa somente de acréscimo, com descarte por rodada & O histórico nunca fica inconsistente, mesmo com interrupções. \\ \hline
7 & Validação local das entradas das ferramentas & Entradas transmitidas em fluxo podem chegar truncadas; nada inválido é executado. \\ \hline
8 & Detecção de fala própria sobre o microfone & Medidor de nível real, cancelamento e pausa de fim de fala menor. \\ \hline
9 & Reconhecimento local e restrito da frase de ativação & Privacidade e exatidão em frases curtas. \\ \hline
10 & SQLite com histórico somente de acréscimo e cadeia de resumos & Memória confiável e auditável em um único arquivo local. \\ \hline
11 & Renderização prévia do avatar com cache invalidado pelo resumo do código & A animação apenas troca imagens; o cache nunca fica desatualizado. \\ \hline
12 & Testes de arquitetura & A regra de dependência e os princípios SOLID não dependem apenas de disciplina. \\ \hline
13 & Despacho por tabelas (intenções, ferramentas e períodos) & Baixa complexidade e extensão sem editar o código existente. \\ \hline
14 & Canvas diferencial e halo opaco & O redesenho de imagens com transparência dominava o custo; enviar só o que muda e evitar transparência reduziu o repouso de 29\% para 5\% de um núcleo. \\ \hline
15 & Taxa de quadros adaptativa e fila de trabalho com orçamento & Fluidez quando há movimento, economia em repouso e nenhum engasgo ao converter imagens. \\ \hline
16 & Corte do silêncio e síntese paralela de frases & As frases se emendam como na fala natural e a primeira palavra sai antes. \\ \hline
17 & Modelo local em camadas, com ferramentas seguras & Um modelo pequeno conversa bem, mas escolhe ferramentas mal; ações destrutivas ficam com comandos determinísticos. \\ \hline
18 & Microfone escolhido pelo servidor de áudio, com variável definida só durante a abertura & Evita disputar o dispositivo com o PipeWire e não afeta os programas abertos pelo assistente. \\ \hline
19 & Estado da inicialização lido dos próprios arquivos & A tela nunca contradiz o sistema, mesmo se o usuário mexer nos atalhos por fora. \\ \hline
20 & Mudança do banco por cópia conferida, sem apagar o anterior & Nenhuma operação sobre os dados do usuário é destrutiva ou irreversível. \\ \hline
\end{tabularx}
\fonte{elaboração própria (2026).}
\end{singlespace}

\chapter{Construção do sistema: passo a passo}

Este capítulo relata, em ordem cronológica, os vinte incrementos que levaram o Robofael da primeira
versão funcional à versão 3.8.0. Para cada incremento são descritos o pedido que o motivou, o que foi
construído, as decisões tomadas, os problemas encontrados e a forma de verificação. A
Tabela~\ref{tab:evolucao} resume a evolução.

\begin{singlespace}
\legendatabela{Evolução do sistema por incremento}\label{tab:evolucao}
\small
\begin{tabularx}{\textwidth}{C{0.8cm}C{1.6cm}LC{2.2cm}}
\toprule
\textbf{Nº} & \textbf{Versão} & \textbf{Incremento} & \textbf{Testes} \\ \midrule
\endfirsthead
\toprule
\textbf{Nº} & \textbf{Versão} & \textbf{Incremento} & \textbf{Testes} \\ \midrule
\endhead
1  & --     & Protótipo: avatar a partir das fotos, voz, agenda e modelo de linguagem & -- \\
2  & --     & Busto centralizado e boca suave & -- \\
3  & --     & Perfil de voz (ajuste de tom e ritmo) & -- \\
4  & --     & Nitidez, camisa futurista e boca natural & -- \\
5  & 1.0.0  & Estrutura profissional, documentação e testes unitários & 98 \\
6  & 2.0.0  & Refatoração completa: arquitetura hexagonal e pipeline em fluxo contínuo & 168 \\
7  & 2.1.0  & Renomeação para Robofael & 168 \\
8  & --     & Gravação da voz do usuário e criação do perfil real & 168 \\
9  & 2.2.0  & Ativação por voz local (“ok Robofael”) & 211 \\
10 & 2.3.0  & Nova frase de ativação (“ok Fael”) com gramática restrita & 247 \\
11 & 3.0.0  & Memória de longo prazo, “o que eu fiz”, janela fixa, configurações e programas & 351 \\
12 & 3.0.1  & Ganho automático do microfone e aviso de voz baixa & 357 \\
13 & 3.1.0  & Internacionalização (português e inglês) & 508 \\
14 & 3.2.0  & Instalador padrão para Ubuntu & 508 \\
15 & 3.3.0  & Revisão SOLID e código limpo, com testes de arquitetura & 529 \\
16 & 3.4.0  & Desempenho e fluidez guiados por medição; tipagem completa & 546 \\
17 & 3.5.0  & Nome técnico \cod{robofael}, com compatibilidade com os nomes antigos & 549 \\
18 & 3.6.0  & Modelo de linguagem local (Ollama) com respostas em camadas & 576 \\
19 & 3.7.0  & Escolha do microfone e abertura junto com o computador & 601 \\
20 & 3.8.0  & Local e backup do banco; modelo em uso e modelos compatíveis & 643 \\
\bottomrule
\end{tabularx}
\fonte{elaboração própria (2026), a partir do registro de mudanças do projeto.}
\end{singlespace}

\section{Incremento 1: protótipo funcional}

\textbf{Pedido.} Criar um programa em Python que manipulasse as imagens do usuário, presentes na pasta
do projeto, para transformá-las no avatar de um assistente pessoal de atividades diárias que ouve e
responde por voz.

\textbf{Construção.} Foram usadas duas fotografias: a do crachá e uma \en{selfie}, ambas já sem
fundo. Em cada uma foram medidas manualmente as posições da boca e do queixo. A animação foi obtida pela
deformação suave da região da mandíbula para baixo, com o preenchimento do espaço aberto por um interior
escuro de boca, em oito níveis de abertura. Um brilho ciano nas bordas deu ao avatar a aparência de
assistente digital. A voz foi implementada com reconhecimento de fala em português pelo serviço do Google
e síntese pela voz neural \cod{pt-BR-AntonioNeural}, com o gTTS como reserva. A abertura da boca passou a
acompanhar o volume real do áudio, quadro a quadro. A agenda de tarefas e lembretes falados era salva em
um arquivo JSON, e o modelo Claude, quando disponível, operava a agenda por meio de ferramentas; sem
chave de acesso, um modo offline atendia a comandos básicos.

\textbf{Verificação.} A interface foi executada com os dois avatares e a fala sincronizada com a
animação. A integração com o modelo de linguagem não pôde ser exercitada de imediato, por falta de
credenciais na máquina, e foi registrada como pendente.

\section{Incremento 2: busto centralizado e boca suave}

\textbf{Pedido.} Exibir apenas um busto, centralizado, e suavizar a animação da boca.

\textbf{Construção.} Mediu-se o topo e o centro da cabeça em cada foto, e o recorte passou a ir de um
pouco acima da cabeça até os ombros, centralizado na cabeça, e não na foto inteira. Com isso, a
\en{selfie}, que estava deslocada para o lado, ficou centralizada. Uma moldura oval com bordas suaves
passou a esconder os cortes das fotos. Para a boca, o número de aberturas subiu de 8 para 16, mais
próximas entre si quando a boca está quase fechada, onde os saltos eram mais visíveis. Cada quadro passou
a ser desenhado no dobro do tamanho e depois reduzido (superamostragem), eliminando o serrilhado. O
volume da voz passou a ser lido 60 vezes por segundo, sem picos isolados nem ruído de fundo. A boca
passou a abrir rápido e fechar devagar, como na fala real, e a animação foi para cerca de 60 quadros por
segundo.

\textbf{Verificação.} Capturas de tela confirmaram o enquadramento e a passagem por aberturas
intermediárias durante a fala.

\section{Incremento 3: perfil de voz}

\textbf{Pedido.} Registrar a voz do usuário e fazer o assistente falar com o mesmo tom. Diante das
alternativas, ajustar apenas o tom de uma voz neural ou clonar o timbre localmente, o usuário escolheu
a primeira.

\textbf{Construção.} Foi criado um gravador de voz com interface própria, que apresenta frases para
leitura, detecta o fim de cada uma e permite regravar. Com pelo menos três gravações, o programa estima o
tom médio por autocorrelação e a velocidade de fala. Em seguida, escolhe a voz neural em português de tom
mais próximo e ajusta tom e ritmo até que a voz sintética fique a menos de 3 Hz da voz do usuário. O
perfil é salvo em arquivo JSON.

\textbf{Verificação.} O medidor de tom acertou sinais de teste com erro máximo de 1 Hz. Em uma
simulação com um “usuário” 25 Hz mais grave e 15\% mais lento, o perfil resultou em um ajuste de
$-22$ Hz, com o tom final a 2 Hz do alvo, e em um ritmo de $-14\%$.

\section{Incremento 4: nitidez, camisa futurista e boca natural}

\textbf{Pedido.} Aumentar a nitidez do avatar, trocar a roupa por uma camisa de corte futurista, tornar
o movimento da boca mais natural e manter o busto centralizado.

\textbf{Construção.} Constatou-se que a \en{selfie} sem fundo tinha apenas 433 por 577 pixels, mas o
arquivo original tinha 1200 por 1600 pixels na mesma proporção. Passou-se então a usar a imagem sem fundo
apenas como máscara de contorno, aplicada sobre a foto em alta resolução. A foto do crachá, disponível
apenas em 500 pixels, recebeu ampliação de alta qualidade, filtro de nitidez e leve aumento de
contraste. A camisa original foi substituída por uma camisa azul-grafite desenhada por cima, com gola
alta metálica de borda luminosa, costuras em V com brilho ciano, placas nos ombros com pequenos LEDs e
um emblema no peito. As dobras do tecido foram preservadas aproveitando-se a luz e a sombra da camisa
original. A boca deixou de ser um corte que descia e passou a ser uma deformação contínua: o lábio
superior sobe levemente, a mandíbula desce, os cantos ficam presos e a abertura tem formato de lente,
com sombra sob os dentes e língua visível nas aberturas maiores. Foram introduzidos três formatos de
vogal (arredondado para “o” e “u”, neutro para “a” e mais largo para “e” e “i”), escolhidos a partir do
timbre do áudio.

\textbf{Verificação.} Prévias geradas com as mesmas imagens e posições da interface, além da execução
da interface com fala simulada, sem erros.

\section{Incremento 5: estrutura profissional e testes (versão 1.0.0)}

\textbf{Pedido.} Organizar o projeto com estrutura profissional, pronta para versionamento, com um
arquivo README que explicasse a arquitetura, a forma de execução e os fluxos, e com testes unitários.

\textbf{Construção.} Os scripts soltos foram reorganizados em um pacote Python instalável, com
subpacotes para lógica (agenda, conversa, ferramentas, modo offline), áudio, avatar e interface, e com a
configuração centralizada e ajustável por variáveis de ambiente. As fotos foram movidas para uma pasta de
recursos com nomes legíveis. A lógica foi separada do microfone, do alto-falante e da tela, e o cliente
do modelo, o relógio e as pastas tornaram-se substituíveis por versões simuladas, o que viabilizou os
testes sem hardware. Os dados pessoais (agenda, perfil e gravações) e as imagens geradas foram colocados
em pastas excluídas do controle de versão. Foram criados um Makefile, o arquivo de projeto
\cod{pyproject.toml} e o README.

\textbf{Verificação.} 98 testes unitários aprovados em cerca de três segundos e análise de estilo sem
ocorrências.

\section{Incremento 6: refatoração completa (versão 2.0.0)}

\textbf{Pedido.} Refatorar o projeto do início ao fim, com nível de engenharia de grandes empresas de
tecnologia, aperfeiçoando a usabilidade, o desempenho e a organização arquitetural.

\textbf{Diagnóstico.} Antes de qualquer alteração, foi feito um diagnóstico. O fluxo era totalmente
serial: esperava-se a resposta inteira do modelo, sintetizava-se todo o áudio e só então a fala começava,
com três a seis segundos de silêncio. O microfone usava uma função da biblioteca que funcionava como
caixa-preta, sem nível de áudio nem cancelamento. Os estados eram cadeias de caracteres alteradas por
várias threads. A renderização do avatar bloqueava a abertura da janela. Os lembretes eram verificados
a cada 15 segundos. A escrita direta do JSON podia corromper o arquivo. Por fim, a classe principal da
interface conhecia todas as tecnologias.

\textbf{Construção.} A solução foi reescrita em arquitetura hexagonal, com:
\begin{enumerate}[label=\alph*)]
  \item domínio tipado com classes de dados no lugar de dicionários;
  \item máquina de estados explícita e barramento de eventos;
  \item pipeline em fluxo contínuo de três estágios, com divisor incremental de frases;
  \item áudio decodificado em memória e envelope da boca calculado antes da reprodução;
  \item cache em disco das falas sintetizadas;
  \item renderização paralela do avatar fora da thread da interface, com cache invalidado pelo resumo do
        código;
  \item inicialização paralela dos serviços em segundo plano;
  \item agendador que dorme até o próximo lembrete;
  \item detecção de fala própria, com medidor de nível real e pausa de fim de fala de 0,7 s;
  \item escrita atômica da agenda e preservação de arquivos corrompidos;
  \item histórico da conversa somente de acréscimo, com descarte por rodada;
  \item validação local das entradas das ferramentas.
\end{enumerate}
Foram adotados o verificador de tipos, a análise estática, os ganchos de \en{pre-commit} e a integração
contínua com servidor gráfico virtual.

\textbf{Problemas corrigidos.} Foram corrigidos cinco problemas: um cancelamento que se perdia quando o
ESPAÇO era pressionado duas vezes rapidamente; uma notificação perdida no agendador; um erro de ponto
flutuante no limiar de tempo da detecção de fala; o acesso à interface do gravador fora da thread
principal; e valores padrão da configuração incompatíveis com classes de dados compactas.

\textbf{Verificação.} 168 testes, incluindo um teste que prova que a primeira frase toca antes de o
modelo terminar de gerar a resposta. A janela passou a aparecer em 0,8 s, e a primeira palavra no modo
offline, em cerca de 1,3 s.

\section{Incremento 7: renomeação para Robofael (versão 2.1.0)}

\textbf{Pedido.} Mudar o nome do assistente de EuRobô para Robofael.

\textbf{Construção.} O nome passou a ficar centralizado em um único campo da configuração, ajustável
por variável de ambiente, e foi aplicado ao título da janela, à saudação, às instruções do modelo, à
ferramenta de gravação, aos registros e à documentação. Optou-se por manter o nome técnico do pacote
(\cod{eurobo}) para não quebrar importações, comandos e caminhos existentes. A renomeação técnica
para \cod{robofael} só ocorreu na versão 3.5.0 (Incremento 17), mantendo os comandos e as variáveis
de ambiente antigos como alternativas.

\textbf{Verificação.} Testes do nome padrão, da troca por variável e da apresentação do modelo como
Robofael.

\section{Incremento 8: gravação da voz real do usuário}

\textbf{Pedido.} Gravar a voz do usuário.

\textbf{Construção e resultado.} A ferramenta de gravação foi aberta, e o usuário gravou as cinco
frases. O tom médio medido foi de 116 Hz. A voz base mais próxima foi a \cod{pt-BR-AntonioNeural}, com
137 Hz. O ajuste de tom de $-21$ Hz levou a voz sintética a 118 Hz, a 2 Hz da voz do usuário, e o ritmo
ficou em $-26\%$. As gravações tiveram entre 3,9 e 5,2 segundos de fala cada uma, com tom entre 95 e
119 Hz, e o perfil usa a mediana das cinco para reduzir o efeito de frases interrogativas.

\textbf{Problemas e diagnóstico.} Durante a primeira execução real, o modelo de linguagem não
conectou. A investigação mostrou que a configuração do terminal do usuário direcionava as chamadas da API
para um servidor local de outro fornecedor, onde o modelo não existia. A mensagem de aviso foi
aprimorada para indicar essa causa, sem alterar a configuração pessoal do usuário. Também foi eliminado
o ruído de mensagens nativas do servidor de áudio JACK que aparecia no terminal a cada abertura do
microfone.

\section{Incremento 9: ativação por voz (versão 2.2.0)}

\textbf{Pedido.} Acionar a captura de voz com a frase “ok Robofael”, em vez de pressionar o ESPAÇO.
Entre o reconhecimento local e o remoto da frase de ativação, o usuário escolheu o local.

\textbf{Construção.} Foi implementada a escuta passiva local com o Vosk e um modelo compacto de
31 MB para português, para que o áudio do ambiente não deixasse o computador. Como o modelo não conhece
a palavra “Robofael” e a transcreve por aproximação (por exemplo, “o que roubou fel”), a detecção foi
feita por semelhança fonética, calibrada com transcrições reais: ativa com a frase sozinha, com a frase
seguida de comando e com a frase dita após outra conversa, mas não com “ok Rafael” nem com conversas
comuns. Quando o reconhecedor remoto erra a frase de ativação, como em “Ok Google Fael, me lembra…”, o
comando é recuperado pelo alinhamento das duas transcrições. Um sinal sonoro de dois tons confirma a
ativação.

\textbf{Problemas corrigidos.} A biblioteca de áudio capturava os sinais de interrupção do sistema, e
o programa não fechava com Ctrl+C nem com SIGTERM. A solução foi desativar os tratadores de sinal da
biblioteca e tratar os sinais no laço da interface.

\textbf{Verificação.} 211 testes; reconhecimento real exercitado com a voz neural do perfil injetada no
lugar do microfone.

\section{Incremento 10: frase “ok Fael” e gramática restrita (versão 2.3.0)}

\textbf{Pedido.} Trocar a frase de ativação para “ok Fael”.

\textbf{Construção.} A palavra “Fael”, curta e fora do vocabulário do modelo local, era simplesmente
descartada: “ok Fael” virava “ok” ou “o que foi”. A transcrição livre foi então substituída por uma
gramática restrita, na qual o reconhecedor só pode escolher entre variantes aceitas de “ok Fael” e
frases distratoras parecidas e frequentes. Com isso, o que não é a frase de ativação é absorvido pelas
distratoras. A frase de ativação passou a ser configurável e independente do nome do assistente. Como a
pausa da vírgula dividia a fala em duas partes, a escuta passou a continuar gravando enquanto houvesse
voz após a frase. O comando também passou a ser limpo da frase quando o reconhecedor remoto a transcreve
de forma inesperada, como em “o que Israel anota tarefa…”, que resulta em “anota tarefa…”.

\textbf{Verificação.} Calibração com cinco vozes sintéticas: zero falsos positivos em 115 frases
parecidas e 8 acertos em 8 tentativas na voz do perfil do usuário; 247 testes.

\section{Incremento 11: memória de longo prazo e outras funções (versão 3.0.0)}

\textbf{Pedido.} Adicionar um banco SQLite como memória de longo prazo, com histórico auditável de
tarefas e lembretes, para que o assistente respondesse o que o usuário fez em uma data ou semana do
passado. O pedido incluía ainda tornar fixo o tamanho da janela, porque a tela cheia quebrava a
renderização; criar um botão de configuração de avatar e voz; permitir abrir programas instalados no
Linux; e garantir a integridade de todos os testes.

\textbf{Construção.} O pedido foi dividido em seis frentes:
\begin{enumerate}[label=\alph*)]
  \item \textbf{memória em SQLite}: agenda, histórico e preferências em um único arquivo; histórico
        somente de acréscimo, com gatilhos que recusam alteração e remoção e com cadeia de resumos SHA-256;
        estado e histórico gravados na mesma transação; esquema versionado; importação automática do
        arquivo JSON anterior, renomeado e preservado; verificação de integridade por comando;
  \item \textbf{“o que eu fiz”}: interpretação de períodos (dezenove formas testadas) e cálculo do
        resultado líquido no período, disponível no modo offline e como ferramenta do modelo;
  \item \textbf{janela de tamanho fixo}: tamanhos mínimo e máximo iguais e redimensionamento desativado;
  \item \textbf{configurações}: diálogo com avatar, voz, tom, ritmo, botão para ouvir, uso do perfil
        gravado e ativação por voz, com cancelamento que desfaz o teste e preferências persistidas;
  \item \textbf{programas}: leitura das entradas \cod{.desktop}, sinônimos, navegador padrão,
        desambiguação e abertura sem interpretador de comandos;
  \item \textbf{testes}: cobertura de todas as frentes, incluindo a detecção de quatro tipos de
        adulteração do histórico.
\end{enumerate}

\textbf{Problemas corrigidos.} Um teste revelou que “anota tarefa abrir conta no banco” abria um
programa, em vez de criar a tarefa; os comandos de agenda passaram a ter prioridade. A análise estática
apontou uma função anônima que usava uma variável já descartada pelo tratamento de exceção. As
variáveis da interface eram destruídas fora da thread principal pelo coletor de lixo e passaram a ser
liberadas explicitamente.

\textbf{Verificação.} 351 testes aprovados em três execuções seguidas e execução real com uma cópia dos
dados do usuário, na qual a migração, a consulta ao histórico e a verificação de integridade
funcionaram.

\section{Incremento 12: ganho automático (versão 3.0.1)}

\textbf{Contexto.} Os primeiros testes com a voz real do usuário mostraram que a frase de ativação não
era reconhecida: o reconhecedor ouvia apenas “ok” e “que”, e o pico do sinal ficava em cerca de 580,
quando o necessário era superior a 1000. A investigação, feita somente por leitura, revelou que o
microfone em uso era o interno do computador, com entrada em 37\%, enquanto um microfone USB conectado
não estava sendo usado.

\textbf{Construção.} Foi implementado um controle automático de ganho que amplifica até 8 vezes a voz
que chega baixa, resistente a picos isolados, e um aviso de voz baixa no terminal e na tela, com
orientação ao usuário. As configurações de áudio do sistema não foram alteradas; os comandos para
ajustá-las foram fornecidos ao usuário.

\textbf{Verificação.} 357 testes; após o ajuste, os registros mostraram nove ativações bem-sucedidas
com a voz real do usuário.

\section{Incremento 13: internacionalização (versão 3.1.0)}

\textbf{Pedido.} Permitir que o assistente responda em inglês quando configurado.

\textbf{Construção.} Todas as mensagens foram reunidas em um catálogo único, com as duas línguas lado a
lado, e um tradutor foi injetado em todos os componentes, com troca ao vivo. O idioma passou a afetar as
respostas, a interface, a voz neural (Antonio em português; Andrew, Brian, Guy ou Jenny em inglês), a
voz de reserva, o reconhecimento de fala e as instruções do modelo, cuja conversa recomeça na troca.
Os comandos passaram a ser entendidos nos dois idiomas, independentemente do idioma das respostas, e os
períodos em inglês foram suportados, com datas numéricas seguindo a convenção de cada idioma.

\textbf{Problemas corrigidos.} O diálogo de configurações podia falhar ao se tornar modal antes de
aparecer na tela e passou a tentar novamente. Além disso, as instruções ao modelo deixaram de atribuir
gênero ao usuário.

\textbf{Verificação.} 508 testes, inclusive um que garante que nenhuma tradução está faltando e que os
campos de formatação coincidem nas duas línguas; execução real em inglês com uma cópia dos dados.

\section{Incremento 14: instalador para Ubuntu (versão 3.2.0)}

\textbf{Pedido.} Criar um instalador padrão para ambientes Ubuntu.

\textbf{Construção.} Optou-se por um pacote \cod{.deb}, o formato nativo da distribuição
\cite{debian}. O pacote contém o código, as fotos, o modelo de ativação e todas as dependências Python
em formato \en{wheel}, o que permite instalar sem internet e de forma determinística. As dependências de
sistema, como Python, Tk, PyAudio e PortAudio, são declaradas para o gerenciador de pacotes. Na
instalação, um ambiente Python isolado é criado em \cod{/opt/robofael} e verificado. O pacote fornece
os comandos \cod{robofael} e \cod{robofael-gravar-voz}, atalhos no menu de aplicativos e um ícone
gerado a partir do próprio avatar. Os dados de cada usuário ficam em \cod{~/.local/share/robofael} e o
cache em \cod{~/.cache/robofael}, conforme as convenções XDG. A remoção apaga o programa, mas preserva
os dados do usuário. A janela passou a declarar a classe \cod{robofael}, para que o ambiente gráfico a
associe ao ícone do menu.

\textbf{Problemas corrigidos.} Uma das dependências não era distribuída como \en{wheel} e passou a ser
construída localmente. Além disso, o ambiente isolado reutilizava uma versão da biblioteca de imagens do
sistema sem suporte à interface gráfica, o que foi resolvido pela instalação forçada das versões
incluídas no pacote.

\textbf{Verificação.} Um roteiro de teste instala o pacote em uma pasta temporária e verifica a
integridade dos arquivos, os atalhos, a criação do ambiente sem internet, os lançadores, os caminhos de
dados, o funcionamento do núcleo e a remoção. A janela foi aberta a partir dessa instalação simulada.

\section{Incremento 15: revisão SOLID e código limpo (versão 3.3.0)}

\textbf{Pedido.} Validar o projeto de ponta a ponta, implementando e verificando os conceitos SOLID e de
código limpo, garantindo que os testes continuassem funcionando.

\textbf{Construção.} Esse incremento é detalhado no Capítulo 7. Em síntese, foram feitas as seguintes
alterações:
\begin{enumerate}[label=\alph*)]
  \item o orquestrador foi dividido em três classes;
  \item as políticas de sessão saíram da raiz de composição e da interface;
  \item três despachos condicionais tornaram-se tabelas extensíveis;
  \item contratos foram criados ou movidos para que todas as dependências apontassem para dentro;
  \item uma violação do princípio de Liskov foi eliminada;
  \item funções longas foram decompostas;
  \item regras de complexidade foram incorporadas à análise estática;
  \item testes de arquitetura passaram a verificar as regras automaticamente.
\end{enumerate}

\textbf{Problema corrigido.} Durante a revisão, um teste de interface passou a falhar de forma
intermitente apenas na suíte completa. A investigação, feita com a captura das pilhas de execução no
momento da falha, revelou que a thread de renderização do avatar ficava bloqueada no destrutor de uma
imagem Tk de um teste anterior: o coletor de lixo do Python, acionado naquela thread, chamava a
biblioteca gráfica fora da thread principal. O defeito era latente e havia sido apenas exposto pela
mudança no momento da coleta. A correção recolhe esses objetos na thread principal ao fim de cada teste
de interface, e a fixture duplicada foi unificada.

\textbf{Verificação.} 529 testes, 85\% de cobertura, nenhuma função com complexidade ciclomática acima
de 10 e imagens do avatar idênticas, pixel a pixel, às anteriores à refatoração.

\section{Incremento 16: desempenho e fluidez guiados por medição (versão 3.4.0)}

\textbf{Pedido.} Refatorar novamente o projeto de ponta a ponta, aperfeiçoando a usabilidade e o
desempenho, com uma organização ainda mais moderna e profissional.

\textbf{Diagnóstico.} Como as revisões anteriores já haviam tratado a estrutura, este incremento começou
por medições do aplicativo real, e não pela leitura do código. Com o assistente parado, depois da
saudação, o processo consumia 29\% de um núcleo de processamento. Um microbenchmark do palco mostrou que
cada quadro levava cerca de 13 ms, perto do limite de 16,6 ms de um quadro a 60 por segundo. Medições
isoladas de cada elemento revelaram a causa: qualquer alteração sob o halo (600 por 600 pixels com
transparência) ou sob o avatar obrigava a biblioteca gráfica a refazer, em software, a mistura de
transparência de toda a área, e girar um único arco custava 8 ms. Além disso, a montagem dos 40 halos
bloqueava a thread da interface por 0,75 s antes de a janela aparecer, e a troca de avatar convertia
36 imagens de uma vez, com um engasgo de cerca de 170 ms. Na fala, a medição do áudio sintetizado
mostrou que cada frase trazia cerca de 0,25 s de silêncio no início e de 0,9 a 1,2 s no fim, o que, na
reprodução frase a frase, se transformava em pausas artificiais.

\textbf{Construção.} Foram realizadas as seguintes alterações:
\begin{enumerate}[label=\alph*)]
  \item \textbf{canvas diferencial}: uma camada entre o palco e a biblioteca gráfica guarda o último valor
        enviado de cada item e só repassa coordenadas, em pixels inteiros, e opções que mudaram;
  \item \textbf{halo opaco}: o halo passou a ser gerado já misturado com a cor de fundo, de modo que a
        biblioteca copia os pixels em vez de misturar transparência; o redesenho sob o halo caiu de
        cerca de 8 ms para 0,1 ms;
  \item \textbf{taxa de quadros adaptativa}: um laço de quadros extraído da janela principal roda a 60
        quadros por segundo com o assistente ouvindo, pensando ou falando, ou durante transições de cor,
        e a 30 em repouso, agendando pelo relógio para não acumular atraso;
  \item \textbf{fila de trabalho com orçamento}: os halos passaram a ser gerados em segundo plano, e a
        conversão de halos e quadros do avatar para a biblioteca gráfica é feita aos poucos, com no
        máximo 4 ms por quadro;
  \item \textbf{corte do silêncio}: uma função de processamento de sinal localiza o trecho com fala e
        mantém margens de 30 ms antes e 150 ms depois, preservando consoantes fracas no início;
  \item \textbf{síntese paralela}: até duas frases são sintetizadas ao mesmo tempo e tocadas na ordem
        original; a decodificação continua em uma única thread, porque o mixer de áudio não é feito
        para uso concorrente;
  \item \textbf{pedido digitado durante a fala}: enviar um texto com o assistente ocupado passou a
        interromper o que ele fazia e atender o pedido novo, em vez de descartá-lo; se vários forem
        enviados, vale o mais recente, e uma trava garante que o pedido seja atendido exatamente uma vez;
  \item \textbf{tipagem completa}: todas as funções passaram a ter assinatura tipada, exigência agora
        verificada pelo mypy.
\end{enumerate}

\textbf{Problemas evitados.} Ao tipar a função de recorte do busto, uma primeira versão passou a usar
coordenadas arredondadas também no cálculo da posição da boca, o que a deslocaria em até meio pixel. A
comparação do resumo SHA-256 de todos os quadros, feita antes e depois, detectou a diferença, e a
correção restaurou as imagens idênticas pixel a pixel. Uma medição inicial do consumo em repouso também
se mostrou contaminada pela saudação falada e foi refeita depois que ela terminava, com a versão anterior
medida nas mesmas condições.

\textbf{Verificação.} 546 testes, incluindo uma regressão de desempenho determinística: um palco em
repouso, desenhado em um canvas falso que conta as chamadas, precisa enviar menos de 80 comandos por
segundo, e as imagens grandes só podem ser redesenhadas poucas vezes por segundo. A cobertura ficou em
85,9\%. No aplicativo real, o consumo em repouso caiu de 29\% para 5\% de um núcleo, a janela passou a
aparecer em 0,6 s (antes, de 1,7 a 2,4 s), a primeira palavra da saudação passou a sair de 0,4 a 0,8 s
mais cedo e a saudação completa ficou cerca de 1,6 s mais curta.

\section{Incremento 17: nome técnico robofael (versão 3.5.0)}

\textbf{Pedido.} Renomear o nome técnico do projeto, que ainda era \cod{eurobo}, para \cod{robofael},
sem quebrar comandos e variáveis que o usuário já usasse.

\textbf{Construção.} O pacote Python, os comandos, as variáveis de ambiente, o lançador do pacote
\cod{.deb} e a documentação passaram a usar o nome \cod{robofael}. Os nomes antigos foram mantidos como
alternativas: os comandos \cod{eurobo} e \cod{eurobo-gravar-voz} continuam instalados como apelidos, e
cada variável \cod{EUROBO_*} ainda é lida quando a \cod{ROBOFAEL_*} correspondente não está definida,
com um aviso de depreciação no registro. O lançador do pacote faz a mesma substituição, para que uma
pasta de dados já configurada pelo nome antigo continue sendo usada.

\textbf{Verificação.} 549 testes, incluindo a leitura dos nomes antigos com aviso e a precedência do
nome novo; um teste de arquitetura garante que o nome antigo só aparece no código de compatibilidade, e
o roteiro do pacote confere que o lançador respeita a variável antiga da pasta de dados.

\section{Incremento 18: modelo de linguagem local (versão 3.6.0)}

\textbf{Pedido.} Permitir que o assistente converse e entenda pedidos em linguagem livre sem o Claude,
com um modelo executado no próprio computador.

\textbf{Construção.} Foi criado um adaptador para o Ollama \cite{ollama}, servidor de modelos locais
acessado por HTTP, escrito apenas com a biblioteca padrão. A raiz de composição passou a escolher o
melhor cérebro disponível, na ordem Claude, modelo local e modo offline, com uma variável que restringe
a escolha. O modelo padrão, de 3 bilhões de parâmetros, foi escolhido por responder em tempo aceitável
no processador, sem placa de vídeo.

\textbf{Problema e decisão.} Os primeiros testes mostraram que um modelo pequeno é bom para conversar,
mas pouco confiável para escolher ferramentas: um pedido de curiosidade abriu o navegador e um
desabafo listou as tarefas. Por isso, a resposta passou a ser montada em camadas:
\begin{enumerate}[label=\alph*)]
  \item comandos que o modo offline reconhece são atendidos por ele, na hora e sempre do mesmo jeito;
  \item pedidos de agenda em linguagem livre vão ao modelo apenas com as ferramentas do assunto e
        somente as seguras (criar e consultar); concluir, remover e abrir programas ficam com os
        comandos explícitos;
  \item o restante é conversa e vai ao modelo sem ferramentas, com o texto em fluxo contínuo;
  \item se o servidor local não responder, a fala é atendida pelo modo offline, e o modelo só volta a
        ser tentado depois de um intervalo.
\end{enumerate}

Chamadas de ferramenta escritas como texto pelo modelo são executadas ou descartadas, mas nunca
faladas. O histórico é compactado a cada troca, e o modelo é carregado na memória antes da primeira
fala, com as mesmas opções da conversa, para evitar uma segunda carga.

\textbf{Verificação.} 576 testes. As camadas, as ferramentas seguras, a queda para o modo offline e o
cliente HTTP são testados contra um servidor falso; no uso real, a primeira palavra de uma conversa
saiu em 1 a 2,5 s, e pedidos com ferramentas levaram de 15 a 45 s no processador.

\section{Incremento 19: escolha do microfone e inicialização (versão 3.7.0)}

\textbf{Pedido.} Permitir selecionar o microfone entre os dispositivos do computador e configurar o
sistema para abrir junto com o computador ou pelo comando \cod{robofael} no terminal.

\textbf{Construção.} Duas portas novas foram criadas, uma para listar os dispositivos de áudio e outra
para a inicialização. Os microfones são listados pelo servidor de áudio (PipeWire \cite{pipewire} ou
PulseAudio), sem os monitores das saídas de som, e a escolha é salva nas preferências. Abrir o
dispositivo físico diretamente disputaria o aparelho com o servidor de áudio e exigiria taxa e canais
exatos; por isso, o microfone continua sendo aberto pelo dispositivo padrão, e a escolha é indicada ao
servidor por variáveis de ambiente definidas apenas durante a abertura, sob uma trava, para não serem
herdadas pelos programas que o assistente abre. A troca vale na próxima escuta: o microfone é
recalibrado e a escuta passiva é reaberta no novo dispositivo. A inicialização usa um atalho de sessão
no padrão de entradas \cod{.desktop} \cite{freedesktop}, com atraso de 5 s e verificação de que o
programa ainda existe, e atalhos em \cod{~/.local/bin} para o terminal. O estado é lido dos próprios
arquivos, e desligar só apaga atalhos criados para este programa. A tela de configurações passou a ter
duas colunas.

\textbf{Verificação.} 601 testes. Uma fixture passou a apontar a pasta pessoal para uma pasta
temporária, para que nenhum teste altere a configuração real do usuário. No aplicativo real, a lista
do servidor de áudio confirmou que a gravação passou a sair do microfone escolhido, e não do padrão.

\section{Incremento 20: banco, backup e modelos de linguagem (versão 3.8.0)}

\textbf{Pedido.} Permitir escolher onde o banco SQLite é salvo e fazer backup dele em qualquer local, e
exibir nas configurações qual modelo de linguagem está em uso e quais modelos compatíveis estão
disponíveis no Claude e no Ollama.

\textbf{Construção.} O banco ganhou as operações de mudar de lugar e de copiar, ambas com o assistente
em uso:
\begin{enumerate}[label=\alph*)]
  \item a cópia usa a interface de backup do SQLite \cite{sqlite}, consistente mesmo com gravações em
        andamento, e é conferida por uma verificação de integridade antes de ser aceita;
  \item ao mudar de lugar, a conexão é trocada sob a mesma trava das gravações, e o arquivo antigo é
        renomeado com a data, nunca apagado; a mudança é recusada se já houver um banco no destino;
  \item como o banco não pode guardar o próprio endereço, o local escolhido fica em um pequeno arquivo
        na pasta de dados; se o banco escolhido não estiver acessível na abertura, como em um disco
        desconectado, o assistente avisa e usa o banco padrão apenas naquela vez, sem criar um banco
        vazio no lugar;
  \item o backup é um arquivo único, sem o diário auxiliar, e informa quantos registros do histórico
        foram copiados.
\end{enumerate}

Para os modelos, cada provedor ganhou um catálogo: o do Claude lista os modelos liberados para a chave,
e o do Ollama lista os modelos instalados capazes de conversar e de chamar ferramentas, com o tamanho.
As consultas são feitas em paralelo, fora da thread da interface, e a falta de chave ou um servidor
fora do ar viram uma explicação na própria lista.

\textbf{Verificação.} 643 testes, incluindo a mudança do banco com gravações antes e depois, a recusa
de sobrescrever outro banco, o banco ausente na abertura, a verificação da cadeia de resumos na cópia e
os catálogos com respostas simuladas. No aplicativo real, a tela mostrou o modelo local em uso e os
dois modelos instalados, e o backup foi conferido.

\section{Lições aprendidas no processo}

O processo evidenciou algumas lições, discutidas novamente no Capítulo 9:

\begin{enumerate}[label=\alph*)]
  \item medições com dados reais, como o volume da voz do usuário e transcrições efetivamente produzidas
        pelos reconhecedores, foram decisivas para resolver problemas que testes com vozes sintéticas
        não revelavam;
  \item decisões com impacto na privacidade ou na experiência do usuário foram submetidas a ele, e
        configurações pessoais do sistema nunca foram alteradas sem autorização;
  \item defeitos de concorrência só apareceram sob carga ou em execuções repetidas, o que justificou
        executar a suíte várias vezes seguidas antes de cada entrega;
  \item otimizar a partir de medições, e não de suposições, levou a ganhos grandes com alterações
        pequenas: o custo dominante estava no redesenho da biblioteca gráfica, e não no código Python;
  \item a arquitetura definida na versão 2.0.0 absorveu sem retrabalho estrutural todas as funções
        seguintes (ativação por voz, memória, programas, idiomas, instalador, modelo local, escolha do
        microfone e mudança do banco), cada uma como um adaptador novo atrás de uma porta, o que confirma
        o valor de investir cedo na separação de responsabilidades;
  \item um modelo de linguagem pequeno precisa de limites explícitos: restringir as ferramentas que ele
        pode usar e deixar as ações destrutivas para comandos determinísticos tornou o modelo local útil
        sem expor o usuário a ações indesejadas;
  \item operações sobre os dados do usuário foram desenhadas para nunca destruir informação: o arquivo
        antigo é renomeado, o destino ocupado é recusado e toda cópia é conferida antes de ser usada.
\end{enumerate}

\chapter{Revisão segundo os princípios SOLID e o código limpo}

Este capítulo detalha o último incremento: a validação do projeto de ponta a ponta à luz dos princípios
SOLID \cite{martin2003} e das práticas de código limpo \cite{martin2009}. A revisão seguiu a disciplina
de refatoração descrita por \textcite{fowler2018}: passos pequenos, comportamento preservado e suíte de
testes executada após cada passo.

\section{Diagnóstico}

O diagnóstico combinou três fontes: a leitura dos módulos maiores; a análise estática com regras
adicionais de complexidade (complexidade ciclomática, número de retornos, de ramificações e de comandos
por função), de nomenclatura e de parâmetros sem uso; e o levantamento automático do grafo de
importações entre as camadas. Os principais achados foram:

\begin{enumerate}[label=\alph*)]
  \item \textbf{SRP}: o orquestrador, com 343 linhas, acumulava o ciclo de conversa, o pipeline de fala
        em três estágios e a escuta passiva; o método que transformava texto em fala tinha complexidade
        ciclomática 20 e 57 comandos. A raiz de composição continha regras do produto (escolha de voz,
        idioma, ativação e saudação), e a janela principal aplicava preferências diretamente nos serviços;
  \item \textbf{OCP}: o executor de ferramentas, o interpretador de intenções do modo offline e o
        interpretador de períodos eram cadeias de condicionais (complexidade 12, 12 e 13), que precisavam
        ser editadas a cada nova ferramenta, intenção ou expressão;
  \item \textbf{DIP}: o núcleo da sessão dependia da classe concreta do banco SQLite; o domínio importava
        contratos da camada de portas (sentido inverso ao desejado); e a ferramenta de gravação de voz
        criava adaptadores por conta própria, sem passar pela raiz de composição;
  \item \textbf{ISP}: a interface acessava o sintetizador e o reprodutor por inspeção dinâmica de
        atributos, sem contrato, e o gravador guardava o microfone como um objeto de tipo indeterminado;
  \item \textbf{LSP}: o repositório JSON legado declarava a assinatura do contrato de persistência da
        agenda, mas ignorava os registros de auditoria, de modo que substituí-lo pelo SQLite perderia o
        histórico em silêncio;
  \item \textbf{código limpo}: uma função de desenho da camisa do avatar com 57 comandos e variáveis de
        uma letra maiúscula; a extração do comando de voz com complexidade 12 e limiares numéricos sem
        nome; um diálogo de configurações com dez parâmetros; e parâmetros não usados em funções anônimas.
\end{enumerate}

\section{Responsabilidade única}

\subsection{Divisão do orquestrador}

O orquestrador foi dividido em três classes, cada uma com um único motivo para mudar:

\begin{enumerate}[label=\alph*)]
  \item \cod{PipelineFala} (\cod{aplicacao/pipeline.py}) transforma um fluxo de texto em fala. Cada
        estágio virou um método curto (pensar, sintetizar e tocar), e o retorno é um objeto imutável com o
        texto falado e a indicação de interrupção;
  \item \cod{VigiaAtivacao} (\cod{aplicacao/vigia.py}) cuida exclusivamente da escuta passiva: reserva
        e libera o microfone sob uma trava, conta falhas seguidas e se desliga com aviso. Ela recebe por
        injeção duas funções: uma que informa se há interação em andamento e outra a ser chamada quando a
        ativação ocorrer;
  \item \cod{Orquestrador} passou a apenas coordenar: decide quando ouvir, pensar e falar, garante uma
        interação por vez e encaixa os anúncios entre as interações.
\end{enumerate}

A interface pública do orquestrador foi preservada. Para isso, atributos como o reprodutor, o
sintetizador e o detector de ativação tornaram-se propriedades que delegam aos colaboradores, e os
testes existentes continuaram válidos sem alteração. O tamanho do módulo caiu de 343 para 255 linhas, e
o método de fala passou de complexidade 20 para 2.

\subsection{Sessão separada da composição e da interface}

As políticas da sessão foram movidas para \cod{aplicacao/sessao.py}: precedência do idioma, escolha da
voz (preferência, perfil gravado e padrão), ativação inicial, texto da saudação, troca de idioma ao vivo
e aplicação das preferências salvas. A raiz de composição ficou restrita a escolher e conectar
implementações. A janela principal, por sua vez, deixou de aplicar regras: ao salvar as configurações,
ela delega à sessão e apenas retraduz os componentes visuais e troca o avatar.

\section{Aberto para extensão, fechado para modificação}

Três despachos condicionais foram substituídos por estruturas extensíveis:

\begin{enumerate}[label=\alph*)]
  \item \textbf{ferramentas do modelo}: o executor mantém uma tabela que associa o nome de cada
        ferramenta a um método. A execução reduziu-se a validar a entrada e chamar a ação correspondente
        (complexidade 1). Acrescentar uma ferramenta significa declarar seu esquema e registrar um método,
        e um teste garante que toda ferramenta declarada tem uma ação e vice-versa;
  \item \textbf{intenções do modo offline}: cada intenção (histórico, anotar, concluir, lembrar, abrir
        programa, listar pendentes, hora e data) é um método que devolve a resposta ou nada. O modelo
        percorre a cadeia em ordem e responde com a primeira intenção que se aplicar, no padrão de cadeia
        de responsabilidade \cite{gamma1994}. A ordem da cadeia expressa a regra de negócio de que comandos
        de agenda têm prioridade sobre a abertura de programas;
  \item \textbf{períodos}: cada forma de expressar um período virou uma função interpretadora pequena,
        que recebe um contexto imutável (texto normalizado, data de referência e idioma) e devolve o
        período ou nada. A lista ordenada, do mais específico para o mais genérico, substituiu a função
        com complexidade 13 por uma de complexidade 1.
\end{enumerate}

\section{Substituição de Liskov}

O banco SQLite foi verificado contra todos os contratos de persistência (agenda, histórico,
preferências e o contrato combinado). O repositório JSON, que desde a versão 3.0.0 serve apenas para
ler o formato legado durante a migração, não guarda o histórico de auditoria. Sua assinatura, porém,
aceitava os registros e os descartava, o que o tornava um falso substituto: o verificador de tipos
aceitaria usá-lo no lugar do banco, e a memória de longo prazo seria perdida em silêncio. O parâmetro foi
removido. Com isso, o repositório deixou de satisfazer o contrato da agenda, e o próprio verificador de
tipos passou a impedir a substituição indevida, aplicando a noção comportamental de subtipo de
\textcite{liskov1994}.

\section{Segregação de interfaces}

Foram criadas duas portas pequenas para eliminar o acesso sem contrato:
\begin{enumerate}[label=\alph*)]
  \item \cod{SintetizadorConfiguravel}: um sintetizador cuja voz, tom e ritmo podem ser lidos e trocados
        ao vivo, usado pela tela de configurações e pela troca de idioma;
  \item \cod{CapturaAudio}: a capacidade de gravar uma frase sem transcrevê-la, usada pela ferramenta de
        gravação da voz.
\end{enumerate}
A porta \cod{PersistenciaAssistente} foi definida como a combinação dos contratos de agenda, histórico e
preferências (na versão 3.8.0, somou-se o de arquivo do banco), usada apenas pelo núcleo da sessão; os demais componentes continuam dependendo somente da
parte de que precisam. A tipagem mais precisa revelou um caso real: o gravador poderia, em tese,
iniciar uma gravação antes de o microfone existir, e o código passou a repassar explicitamente o
microfone já verificado.

\section{Inversão de dependência}

Três ajustes alinharam todas as dependências à regra de dependência:

\begin{enumerate}[label=\alph*)]
  \item o contrato de persistência exigido pela agenda e a entidade que representa um programa
        instalado foram movidos para o domínio, que passou a não importar nada além do catálogo de
        mensagens; a camada de portas os reexporta, mantendo a compatibilidade;
  \item o núcleo da sessão passou a depender da porta de persistência, e não da classe do SQLite; a
        leitura do perfil de voz passou a ser injetada como uma função;
  \item a ferramenta de gravação de voz passou a obter o microfone e a reprodução de amostras por meio
        de fábricas da raiz de composição, deixando de importar adaptadores diretamente.
\end{enumerate}

\section{Práticas de código limpo}

Além das refatorações associadas aos princípios, foram aplicadas as seguintes práticas:

\begin{enumerate}[label=\alph*)]
  \item \textbf{funções pequenas}: o desenho da camisa do avatar foi decomposto em etapas com nome
        (sombreado, tecido, placas dos ombros, gola metálica e linhas luminosas), com um objeto de valor
        para a geometria da gola; a extração do comando de voz foi dividida em três estratégias nomeadas;
  \item \textbf{nomes que revelam intenção}: variáveis de uma letra maiúscula foram renomeadas para
        largura e altura, e os limiares numéricos da extração do comando receberam constantes nomeadas;
  \item \textbf{poucos argumentos}: o diálogo de configurações passou a receber dois objetos-parâmetro,
        um com as opções oferecidas e outro com as ações dos botões, em vez de dez argumentos;
  \item \textbf{sem parâmetros mortos}: parâmetros não usados foram removidos ou marcados explicitamente;
  \item \textbf{sem duplicação}: a preparação da janela de teste, repetida em dois arquivos de teste,
        foi unificada.
\end{enumerate}

Para garantir que o comportamento visual não mudasse, o resumo SHA-256 de todos os 72 quadros dos dois
avatares foi calculado antes e depois da decomposição da função de desenho, e os resumos coincidiram.

\section{Regras permanentes}

Para que a qualidade obtida não se perdesse nas evoluções seguintes, as regras passaram a ser
verificadas automaticamente em toda execução de \cod{make verificar} e na integração contínua:

\begin{enumerate}[label=\alph*)]
  \item complexidade ciclomática máxima de 10 por função \cite{mccabe1976};
  \item limites de retornos, ramificações e comandos por função;
  \item convenções de nomenclatura e ausência de parâmetros sem uso;
  \item verificação estática de tipos sem ocorrências;
  \item cobertura mínima de 80\% das linhas e ramificações;
  \item testes de arquitetura, descritos na Seção~\ref{sec:testes-arquitetura}.
\end{enumerate}

As exceções às regras foram documentadas na própria configuração com sua justificativa. Uma delas
permite nomes de exceção do domínio em português, como \cod{TarefaNaoEncontrada}, sem o sufixo
\cod{Error}. Outra libera os dublês de teste de usar todos os argumentos dos contratos que imitam.

\section{Métricas antes e depois}

A Tabela~\ref{tab:metricas} compara as funções e os módulos alterados.

\begin{singlespace}
\legendatabela{Métricas antes e depois da revisão}\label{tab:metricas}
\small
\begin{tabularx}{\textwidth}{LC{2.6cm}C{2.6cm}}
\toprule
\textbf{Elemento (métrica)} & \textbf{Antes} & \textbf{Depois} \\ \midrule
\endfirsthead
\toprule
\textbf{Elemento (métrica)} & \textbf{Antes} & \textbf{Depois} \\ \midrule
\endhead
Orquestrador (linhas do módulo) & 343 & 255 \\
Método de fala do orquestrador (complexidade) & 20 & 2 \\
Execução de ferramentas (complexidade) & 12 & 1 \\
Interpretação offline (complexidade) & 12 & 3 \\
Interpretação de períodos (complexidade) & 13 & 1 \\
Extração do comando de voz (complexidade) & 12 & 5 \\
Desenho da camisa (comandos na função) & 57 & 22 \\
Diálogo de configurações (parâmetros) & 10 & 5 \\
Maior complexidade ciclomática do projeto & 20 & 10 \\
Ocorrências das novas regras (complexidade, tamanho, nomes e parâmetros) & 22 & 0 \\
Dependências entre camadas fora da regra & 2 & 0 \\
Testes automatizados & 508 & 529 \\
\bottomrule
\end{tabularx}
\fonte{elaboração própria (2026), com dados do ruff e do pytest.}
\end{singlespace}

A redução da complexidade não foi obtida à custa de mais código total: os colaboradores extraídos
somam linhas, mas cada função passou a ser curta, nomeada e testável isoladamente, o que é o objetivo
das práticas de código limpo.

\chapter{Testes e validação}

Este capítulo apresenta a estratégia de testes, as suítes que compõem a validação automatizada, os
testes de arquitetura que verificam os princípios de projeto, o processo de integração contínua e os
procedimentos de validação manual.

\section{Estratégia}

A estratégia segue a pirâmide de testes \cite{cohn2009}. A base é formada por testes unitários rápidos
das regras de domínio, das políticas de aplicação e do processamento de sinais. O nível intermediário
reúne testes de integração dos adaptadores com dublês das bibliotecas externas e do banco SQLite real
em arquivo temporário. No topo ficam poucos testes que abrem a janela Tk real, marcados como lentos e
executados com servidor gráfico virtual na integração contínua.

Três decisões de projeto tornaram essa estratégia viável. A primeira é a arquitetura hexagonal: todo
serviço externo (modelo de linguagem, reconhecedores, sintetizador, reprodutor, lançador de programas)
é acessado por uma porta e pode ser substituído por um dublê \cite{meszaros2007}. A segunda é a injeção
do relógio, que torna determinísticos os testes de lembretes e de períodos. A terceira é a
reprodução simulada com duração configurável, que permite provar propriedades temporais do pipeline,
como a de que a primeira frase começa a tocar antes de o modelo terminar a resposta.

\section{Dublês de teste}

O módulo comum de testes oferece dublês para cada porta:
\begin{enumerate}[label=\alph*)]
  \item modelos roteirizados, que devolvem pedaços de texto com atraso controlado;
  \item um cliente falso do modelo Claude, que reproduz o protocolo de transmissão em fluxo contínuo com
        blocos de texto e de ferramenta;
  \item sintetizadores que registram os textos e podem falhar sob demanda;
  \item reprodutores que registram as faixas tocadas e acionam ganchos durante a reprodução;
  \item reconhecedores com falas roteirizadas e transcrições opcionais;
  \item detectores de ativação que devolvem ativações ou falhas roteirizadas e depois aguardam o
        cancelamento;
  \item lançadores de programas vazios ou com listas controladas;
  \item um servidor HTTP falso do Ollama e catálogos de modelos roteirizados, que respondem com listas
        ou com falhas (sem chave, sem conexão, servidor fora do ar);
  \item listas de microfones e estados de inicialização controlados, sem depender do hardware.
\end{enumerate}
Uma fixture automática aponta a pasta pessoal e a pasta de configuração para diretórios temporários
em todos os testes, de modo que nenhum deles possa criar atalhos de inicialização ou comandos na pasta
real do usuário.
Também são oferecidos geradores de áudio sintético, como senoides em frequências conhecidas e arquivos
WAV temporários, para validar a análise de tom e de envelope.

\section{Suítes de teste}

O Quadro~\ref{qua:suites} descreve as suítes e o número de casos de teste de cada uma, contando as
variações parametrizadas.

\begin{singlespace}
\legendaquadro{Suítes de teste automatizado}\label{qua:suites}
\small
\begin{tabularx}{\textwidth}{|P{4.2cm}|C{1.3cm}|L|}
\hline
\textbf{Suíte} & \textbf{Casos} & \textbf{O que garante} \\ \hline
\endfirsthead
\hline
\textbf{Suíte} & \textbf{Casos} & \textbf{O que garante} \\ \hline
\endhead
\cod{test_i18n} & 190 & Catálogo completo e com os mesmos campos nas duas línguas; comandos, períodos e respostas em inglês; troca ao vivo; precedência da configuração \\ \hline
\cod{test_ativacao} & 79 & Gramática de “ok Fael” e busca fonética com saídas reais do reconhecedor local; frases parecidas que não podem disparar; extração do comando com saídas reais do reconhecedor remoto; vigia liberando o microfone; desligamento após falhas \\ \hline
\cod{test_memoria} & 41 & Dezenove expressões de período, intenção “o que eu fiz”, resultado líquido e respostas \\ \hline
\cod{test_programas} & 29 & Leitura de entradas \cod{.desktop}, extração do pedido, ranqueamento, sinônimos, ambiguidade, abertura sem interpretador de comandos e prioridade da agenda \\ \hline
\cod{test_ollama} & 24 & Respostas em camadas, ferramentas seguras, chamadas escritas como texto, queda para o modo offline, interrupção, histórico, idioma, cliente HTTP contra servidor falso e escolha do cérebro \\ \hline
\cod{test_sqlite} & 18 & Esquema versionado, transação conjunta de estado e histórico, gatilhos de somente acréscimo, detecção de quatro tipos de adulteração, consulta por período, concorrência, migração e preferências \\ \hline
\cod{test_agenda} & 18 & Regras da agenda e transações (persiste e notifica somente em caso de sucesso) \\ \hline
\cod{test_analise_audio} & 23 & Estimativa de tom com erro de até 2\%, envelope de volume e timbre e corte do silêncio das pontas (sem cortar consoantes fracas) \\ \hline
\cod{test_orquestrador} & 17 & Prova do pipeline, interrupção, modo contínuo, anúncios em fila, falhas de modelo, síntese e microfone, corrida do cancelamento e pedidos digitados durante a fala ou a escuta \\ \hline
\cod{test_modelo_offline} & 15 & Comandos, plural, lembretes relativos e despedida \\ \hline
\cod{test_animacao} & 15 & Suavização independente da taxa de quadros e seleção do quadro da boca \\ \hline
\cod{test_modelo_claude} & 14 & Transmissão real do texto, laço de ferramentas, limite de tamanho, recusa, erros da API, descarte da rodada e limite do histórico \\ \hline
\cod{test_arquitetura} & 15 & Regra de dependência entre camadas e princípios SOLID (Seção~\ref{sec:testes-arquitetura}) \\ \hline
\cod{test_configuracoes} & 15 & Conversões de tom e ritmo, precedência das preferências, janela de tamanho fixo e diálogo de configurações, inclusive microfone, inicialização, modelo em uso, mudança e backup do banco \\ \hline
\cod{test_microfone_inicializacao} & 14 & Lista de microfones sem monitores, variável de escolha restrita à abertura, rótulos, reabertura da escuta passiva, atalho de sessão e comandos no terminal sem apagar arquivos de outra origem \\ \hline
\cod{test_banco_e_modelos} & 13 & Backup conferido, mudança do banco com gravações antes e depois, recusa de sobrescrever, banco ausente na abertura, catálogos do Claude e do Ollama e texto da lista de modelos \\ \hline
\cod{test_perfil_voz} & 12 & Criação, leitura e tolerância a arquivos inexistentes ou corrompidos \\ \hline
\cod{test_ferramentas} & 9 & Esquemas e validação de entradas (tipos, campos extras, booleano diferente de inteiro) \\ \hline
\cod{test_avatar} & 9 & Busto centralizado, camisa, tamanho uniforme dos quadros e cache \\ \hline
\cod{test_vad}, \cod{test_estados_barramento}, \cod{test_frases} & 23 & Detecção de fala, transições de estado, entrega de eventos entre threads e divisão de frases \\ \hline
\cod{test_sessao_pipeline_vigia} & 8 & Pipeline e vigia isolados, síntese paralela com ordem preservada e aplicação das preferências \\ \hline
\cod{test_renderizacao} & 8 & Canvas diferencial, fila de trabalho com orçamento, laço de quadros adaptativo, halos opacos e regressão de desempenho do palco em repouso \\ \hline
\cod{test_ganho}, \cod{test_configuracao_tema}, \cod{test_adaptadores_voz} & 21 & Ganho automático, configuração (inclusive os nomes antigos das variáveis) e tema, cache LRU de voz e reprodutor com mixer real \\ \hline
\cod{test_agendador}, \cod{test_repositorio_json}, \cod{test_bootstrap_ui} & 13 & Lembrete no horário exato, formato legado, composição, queda para o modo offline e janela Tk real \\ \hline
\textbf{Total} & \textbf{643} & \\ \hline
\end{tabularx}
\fonte{elaboração própria (2026), a partir da coleta do pytest.}
\end{singlespace}

\section{Testes de arquitetura}\label{sec:testes-arquitetura}

Inspirados nas funções de aptidão arquitetural de \textcite{ford2017}, os testes de arquitetura
transformam as regras de projeto em verificações executáveis. Eles estão no arquivo
\cod{tests/test_arquitetura.py} e verificam:

\begin{enumerate}[label=\alph*)]
  \item \textbf{regra de dependência}: a árvore sintática de cada módulo é analisada, todas as
        importações internas são extraídas e cada uma é comparada com uma tabela que declara, para cada
        camada, as camadas que ela pode importar; qualquer violação é relatada com arquivo, linha e
        sentido da dependência;
  \item \textbf{cobertura das regras}: toda camada existente no pacote precisa ter uma regra declarada,
        para que uma camada nova não escape da verificação;
  \item \textbf{pureza do núcleo}: o domínio e as portas não podem importar bibliotecas de
        infraestrutura (banco de dados, interface gráfica, processamento numérico, áudio, rede, modelo de
        linguagem nem execução de processos);
  \item \textbf{LSP}: o banco SQLite precisa oferecer todos os métodos de todos os contratos de
        persistência, e o repositório legado não pode ter a assinatura que promete guardar o histórico;
  \item \textbf{DIP}: a agenda precisa funcionar com um repositório em memória escrito no próprio teste,
        sem banco de dados;
  \item \textbf{OCP}: toda ferramenta declarada ao modelo precisa ter uma ação registrada e vice-versa, e
        uma intenção nova acrescentada à cadeia do modo offline precisa funcionar sem alterar o código
        existente;
  \item \textbf{SRP}: os métodos do pipeline e da vigia não podem voltar ao orquestrador, e as políticas
        da sessão não podem voltar à raiz de composição.
\end{enumerate}

\section{Análise estática e integração contínua}

O comando \cod{make verificar} executa, em sequência, a análise estática (ruff), a verificação de tipos
(mypy) e a suíte completa (pytest). O mesmo conjunto é executado pelos ganchos de \en{pre-commit} e
pela integração contínua no GitHub Actions, em Python 3.10 e 3.12, com o servidor gráfico virtual Xvfb
para os testes de interface. A medição de cobertura falha se o total ficar abaixo de 80\%.

\section{Cobertura}

A cobertura combinada de linhas e ramificações foi de 86,6\% sobre 4.694 comandos. A
Tabela~\ref{tab:cobertura} apresenta a cobertura por camada.

\begin{singlespace}
\legendatabela{Cobertura de testes por camada (linhas e ramificações)}\label{tab:cobertura}
\small
\begin{tabularx}{\textwidth}{LC{2.6cm}C{3.2cm}}
\toprule
\textbf{Camada} & \textbf{Comandos} & \textbf{Cobertura (\%)} \\ \midrule
\endfirsthead
\toprule
\textbf{Camada} & \textbf{Comandos} & \textbf{Cobertura (\%)} \\ \midrule
\endhead
Portas & 85 & 100,0 \\
Internacionalização & 68 & 98,7 \\
Domínio & 551 & 98,2 \\
Aplicação & 758 & 95,5 \\
Avatar & 290 & 95,2 \\
Configuração & 117 & 93,8 \\
Áudio & 260 & 90,8 \\
Adaptadores & 1.297 & 82,1 \\
Apresentação & 1.130 & 76,0 \\
Raiz de composição & 136 & 75,0 \\ \midrule
\textbf{Total} & \textbf{4.694} & \textbf{86,6} \\
\bottomrule
\end{tabularx}
\fonte{elaboração própria (2026), com dados do pytest-cov.}
\end{singlespace}

A distribuição é coerente com a arquitetura: as camadas internas, que concentram as regras, têm
cobertura próxima de 100\%. As menores coberturas estão nos adaptadores de hardware (microfone,
áudio nativo e escuta local), cuja execução exige dispositivos reais, e na interface gráfica, cujos
caminhos de desenho são parcialmente exercitados pelos testes com a janela real.

\section{Validação manual}

Os aspectos que dependem de hardware ou de percepção humana foram validados manualmente:
\begin{enumerate}[label=\alph*)]
  \item a aparência do avatar, por meio de prévias geradas com as mesmas imagens e posições da interface;
  \item o reconhecimento da frase de ativação com a voz real do usuário, acompanhado pelos registros
        detalhados, que mostram o texto reconhecido, o volume de pico e o ganho aplicado;
  \item a gravação real da voz e a criação do perfil;
  \item a migração dos dados reais, feita primeiro em uma cópia;
  \item a execução em inglês;
  \item a abertura da janela a partir da instalação simulada do pacote, com a verificação da classe da
        janela e da criação das pastas de dados e de cache do usuário;
  \item a escolha do microfone, conferida na lista de gravações ativas do servidor de áudio, que mostrou
        o assistente gravando do microfone escolhido, e não do padrão;
  \item a tela de configurações com os serviços reais e uma pasta de dados temporária, mostrando o
        modelo local em uso, os modelos instalados e um backup conferido;
  \item o atalho de sessão gerado, validado pela ferramenta \cod{desktop-file-validate}.
\end{enumerate}

\section{Estabilidade da suíte}

Como o sistema é concorrente, uma suíte aprovada uma única vez não basta. Antes de cada entrega, a suíte
completa foi executada repetidamente, inclusive com a máquina sob carga. Esse procedimento revelou e
permitiu corrigir quatro defeitos intermitentes ao longo do projeto: dois testes que dependiam de
tempo sob carga, a liberação de variáveis da interface fora da thread principal e o bloqueio da thread
de renderização pelo coletor de lixo, descrito no Capítulo 6.

\chapter{Resultados e discussão}

Este capítulo apresenta os resultados obtidos em relação aos requisitos e aos objetivos, discute as
decisões de projeto à luz desses resultados e expõe as limitações identificadas.

\section{Atendimento aos requisitos funcionais}

Os 42 requisitos funcionais especificados no Capítulo 4 foram implementados. Destes, 41 são
verificados por testes automatizados. O requisito de instalação (RF36) é verificado por um roteiro
próprio, que instala o pacote em uma pasta temporária e confere seu funcionamento sem alterar o
sistema. Algumas funções dependem de serviços externos ou de hardware, e sua verificação automatizada
usa dublês: a conversa livre com o modelo de linguagem (RF24), o reconhecimento remoto (RF03 e RF05) e
a voz neural (RF06). A escolha do microfone (RF39) foi conferida também no servidor de áudio real.
Nessas funções, a validação em uso real complementou os testes. A rastreabilidade
completa está no Apêndice A.

\section{Atendimento aos requisitos não funcionais}

A Tabela~\ref{tab:desempenho} compara as medições de desempenho da arquitetura inicial, serial, com as
da arquitetura em pipeline introduzida na versão 2.0.0 e mantida até a versão final.

\begin{singlespace}
\legendatabela{Medições de desempenho}\label{tab:desempenho}
\small
\begin{tabularx}{\textwidth}{LP{4.2cm}P{4.2cm}}
\toprule
\textbf{Medida} & \textbf{Arquitetura inicial} & \textbf{Arquitetura final} \\ \midrule
\endfirsthead
\toprule
\textbf{Medida} & \textbf{Arquitetura inicial} & \textbf{Arquitetura final} \\ \midrule
\endhead
Janela visível após abrir & 0,5 a 4 s & 0,8 s \\
Primeira palavra (modo offline, texto) & resposta e voz inteiras & cerca de 1,3 s \\
Primeira palavra (modelo de linguagem) & geração e síntese completas & primeira frase gerada e sintetizada \\
Tempo para interromper & apenas durante a reprodução & cerca de 10 ms, em qualquer estágio \\
Renderizar o avatar sem cache (crachá / \en{selfie}) & 2,9 s / 4,3 s & 1,8 s / 2,2 s, em segundo plano \\
Ler o avatar do cache & cerca de 0,4 s & 0,08 s / 0,2 s \\
Repetir uma fala já sintetizada & nova síntese & 0,1 ms (cache LRU) \\
Atraso de um lembrete & até 15 s & menos de 0,1 s \\
Pausa de fim de fala & 0,9 s fixos & 0,7 s, ignorando ruídos curtos \\
\bottomrule
\end{tabularx}
\fonte{elaboração própria (2026).}
\end{singlespace}

O Incremento 16 (versão 3.4.0) acrescentou um refinamento guiado por medições do aplicativo real,
apresentado na Tabela~\ref{tab:desempenho34}. O consumo em repouso foi medido em uma janela de 20 a 30 s
depois do fim da saudação falada, com as duas versões nas mesmas condições e com o cache do avatar já
pronto.

\begin{singlespace}
\legendatabela{Medições de desempenho e fluidez: versões 3.3.0 e 3.4.0}\label{tab:desempenho34}
\small
\begin{tabularx}{\textwidth}{LP{3.6cm}P{3.6cm}}
\toprule
\textbf{Medida} & \textbf{Versão 3.3.0} & \textbf{Versão 3.4.0} \\ \midrule
\endfirsthead
\toprule
\textbf{Medida} & \textbf{Versão 3.3.0} & \textbf{Versão 3.4.0} \\ \midrule
\endhead
Processamento em repouso & 29\% de um núcleo & 5\% de um núcleo \\
Custo de um quadro do palco (microbenchmark) & 13 ms & 5 ms \\
Janela visível após abrir & 1,7 a 2,4 s & 0,6 a 0,7 s \\
Primeira palavra da saudação & 3,1 a 3,2 s & 2,4 a 2,8 s \\
Duração da saudação & cerca de 10,1 s & cerca de 8,5 s \\
Silêncio entre frases faladas & cerca de 1,2 s & 0,15 s \\
Engasgo ao trocar de avatar & cerca de 170 ms & nenhum (no máximo 4 ms por quadro) \\
Memória residente & cerca de 369 MB & cerca de 368 MB \\
\bottomrule
\end{tabularx}
\fonte{elaboração própria (2026).}
\end{singlespace}

A memória não diminuiu: cerca de 150 MB correspondem às bibliotecas (o SDK do modelo de linguagem,
o reconhecedor local com seu modelo, o mixer de áudio e o processamento numérico) e cerca de 100 MB às
imagens já convertidas para a biblioteca gráfica. Um teste com a devolução explícita de memória ao
sistema operacional recuperou apenas 10 MB, o que indicou uso legítimo, e não fragmentação.

Os requisitos de desempenho (RNF01 a RNF05, RNF26 e RNF27) foram atendidos conforme as tabelas
anteriores. Os requisitos de confiabilidade (RNF06 a RNF09) foram atendidos e são verificados por testes que
simulam falhas de transação, adulterações do histórico, interrupções durante o fluxo de texto, recusas
e erros da API, ausência de microfone e falhas repetidas da escuta passiva. O requisito de exatidão da
ativação (RNF13) foi atendido no conjunto de calibração, com zero falsos positivos em 115 frases
parecidas, e confirmado em uso real com a voz do usuário depois da introdução do ganho automático.
Os requisitos de manutenibilidade (RNF19 a RNF22) foram atendidos e passaram a ser verificados
continuamente: nenhuma dependência fora da regra entre camadas, complexidade máxima de 10, tipos
verificados sem ocorrências, com todas as assinaturas tipadas, e cobertura de 85,9\%, acima do mínimo
de 80\%.

\section{Discussão}

\subsection{Latência percebida}

O maior ganho de experiência veio de uma decisão arquitetural, e não de otimizações locais. Ao
sobrepor a geração de texto, a síntese e a reprodução, o tempo até a primeira palavra deixou de
depender do tamanho da resposta. Esse resultado está de acordo com a observação de
\textcite{nielsen1993} de que respostas em torno de um segundo preservam o fluxo da interação. O custo
dessa escolha é a concorrência, que exigiu mecanismos explícitos: token de cancelamento por interação,
máquina de estados e barramento de eventos.

\subsection{Privacidade e exatidão da ativação}

O uso de um reconhecedor local com gramática restrita mostrou-se superior à transcrição livre para uma
frase de ativação curta e fora do vocabulário. As frases distratoras funcionam como uma classe de
rejeição explícita, o que explica a ausência de falsos positivos na calibração. Os testes com a voz
real, por outro lado, mostraram que a qualidade da captação é tão importante quanto o algoritmo: com o
microfone a 37\% de entrada, nenhum ajuste de gramática resolveria o problema. O ganho automático
melhorou a situação, mas não substitui um volume de entrada adequado.

\subsection{Memória auditável}

O encadeamento de resumos criptográficos, somado aos gatilhos de somente acréscimo, oferece duas
camadas de proteção: o próprio banco recusa alterações feitas pelo sistema e por ferramentas comuns, e a
verificação da cadeia detecta modificações feitas por fora, inclusive com o arquivo aberto em outro
programa. Como os registros guardam a descrição no momento do fato, a pergunta “o que eu fiz” é
respondida corretamente mesmo depois da remoção dos itens.

\subsection{Princípios SOLID na prática}

A revisão final mostrou que uma arquitetura em camadas bem definida não impede, por si só, a erosão
local da qualidade. Mesmo com a separação hexagonal desde a versão 2.0.0, as funções acumularam
condicionais, o orquestrador absorveu responsabilidades e duas dependências passaram a apontar no
sentido errado ao longo de nove incrementos. Os princípios SOLID funcionaram como critérios objetivos
para localizar e corrigir esses pontos, e as funções de aptidão de \textcite{ford2017} tornaram a
correção permanente. Em vez de depender da atenção de quem altera o código, a violação passa a ser
apontada pelo próprio processo de verificação.

Um efeito colateral positivo foi a melhoria da tipagem. Ao substituir acessos dinâmicos por contratos
explícitos, o verificador de tipos encontrou um caso real de uso possivelmente nulo, e a correção do
princípio de Liskov transformou um risco silencioso de perda de dados em um erro de tipo.

\subsection{Desenvolvimento assistido por inteligência artificial}

A experiência indica que um agente de programação acelera substancialmente a implementação, inclusive
de refatorações amplas, mas desloca a responsabilidade pela qualidade para as verificações automáticas
e para a validação com o usuário. Os defeitos mais difíceis do projeto foram identificados por medições
reais ou por execuções repetidas, como o volume baixo do microfone, a configuração que redirecionava a
API e o bloqueio da thread pelo coletor de lixo, e não pela leitura do código. Isso reforça o valor das
práticas de engenharia descritas neste trabalho, que se mostraram complementares ao uso do agente.

\section{Limitações}

As principais limitações identificadas são:
\begin{enumerate}[label=\alph*)]
  \item o reconhecimento de comandos e a voz neural dependem de serviços na internet; sem conexão, a
        conversa livre pode usar o modelo local, mas a voz de reserva também fica indisponível;
  \item o modelo local executado no processador leva de 15 a 45 s em pedidos que usam ferramentas e,
        por ser pequeno, só recebe ferramentas que criam ou consultam: concluir, remover e abrir
        programas exigem os comandos explícitos;
  \item a tela de configurações mostra os modelos disponíveis, mas a troca do modelo ainda é feita por
        variáveis de ambiente e vale ao reabrir o assistente;
  \item a restauração de um backup é manual, com o assistente fechado;
  \item a gramática restrita foi calibrada apenas para a frase “ok Fael”; outras frases usam a busca
        fonética livre, menos precisa;
  \item a calibração da ativação usou cinco vozes sintéticas e a voz de um único usuário real; uma
        avaliação com mais pessoas e ambientes ruidosos daria estimativas mais robustas de exatidão;
  \item os lembretes só são anunciados com o programa aberto; um lembrete vencido com o programa fechado
        é anunciado na abertura seguinte;
  \item o pacote de instalação foi construído e testado para Ubuntu 24.04 com Python 3.12; outras
        versões exigem reconstruí-lo na própria máquina, e as fotografias pessoais seguem dentro do
        pacote;
  \item as posições do rosto em cada fotografia foram medidas manualmente; um novo avatar exige
        cadastrar essas posições no catálogo.
\end{enumerate}

\chapter{Conclusão}

Este trabalho apresentou o Robofael, um assistente pessoal de atividades diárias para Linux que
conversa por voz, é representado por um avatar animado criado a partir das fotografias do usuário,
organiza tarefas e lembretes, mantém uma memória de longo prazo auditável, abre programas instalados,
funciona em português e em inglês e pode ser instalado no Ubuntu por um pacote padrão. Sem o Claude, um
modelo de linguagem local assume a conversa; e a tela de configurações permite escolher o microfone,
abrir o assistente com o computador, mudar o banco de lugar e fazer backups conferidos. O problema de
pesquisa, que buscava construir um assistente responsivo, respeitoso com a privacidade, com histórico
confiável e estrutura interna fácil de evoluir, foi respondido com uma solução que combina decisões
arquiteturais e verificações automáticas.

Os objetivos específicos foram atingidos. O avatar em formato de busto, com boca sincronizada em
36 combinações de abertura e formato, foi gerado a partir das fotografias. O ciclo de conversa em
pipeline reduziu a espera à primeira frase da resposta. A ativação pela frase “ok Fael” é processada
localmente, com gramática restrita e zero falsos positivos na calibração. A memória de longo prazo em
SQLite é somente de acréscimo e protegida por cadeia de resumos SHA-256, e responde perguntas sobre dias,
semanas e outros períodos. A integração com o ambiente Linux abre programas por nome ou sinônimo. A
internacionalização permite trocar o idioma ao vivo. O instalador para Ubuntu foi testado em ambiente
isolado. Por fim, a revisão segundo os princípios SOLID e as práticas de código limpo foi acompanhada de
testes de arquitetura que tornam essas regras permanentes. Os requisitos funcionais e não funcionais
foram especificados, rastreados até os testes e documentados, assim como todo o processo de construção.

Como contribuição, além do próprio artefato de software, o trabalho mostra de forma concreta como os
princípios SOLID se materializam em um sistema concorrente, com integração a serviços externos e
interface gráfica: classes com uma única responsabilidade no lugar de um coordenador sobrecarregado;
tabelas extensíveis no lugar de cadeias de condicionais; contratos pequenos e pertencentes à camada que
deles precisa; e um caso real em que respeitar o princípio de Liskov transformou uma perda silenciosa
de dados em um erro detectado pelo verificador de tipos. O trabalho também mostra que testes de
arquitetura, limites de complexidade e cobertura mínima, verificados continuamente, preservam a
qualidade estrutural ao longo da evolução, inclusive quando a implementação é acelerada por um agente
de inteligência artificial. O Incremento 16 ilustra o valor de otimizar a partir de
medições: com uma camada que só repassa à biblioteca gráfica o que mudou, imagens opacas e uma taxa de
quadros adaptativa, o consumo em repouso caiu de 29\% para 5\% de um núcleo, sem alterar a aparência,
e um teste determinístico passou a impedir que esse ganho se perca. Por fim, os quatro últimos
incrementos mostram que a arquitetura absorveu um novo cérebro, novos dispositivos e novas operações
sobre os dados apenas com portas e adaptadores novos, e que um modelo pequeno se torna útil quando o
sistema limita o que ele pode fazer.

\section{Trabalhos futuros}

A partir das limitações identificadas, sugerem-se as seguintes evoluções:
\begin{enumerate}[label=\alph*)]
  \item adotar reconhecimento de comandos e síntese de voz locais, que, somados ao modelo de linguagem
        local já existente, eliminariam a dependência de internet para as funções principais;
  \item permitir escolher o modelo de linguagem na própria tela de configurações, com troca ao vivo;
  \item restaurar backups pela interface e criar backups automáticos periódicos;
  \item gerar automaticamente a gramática restrita de qualquer frase de ativação, com frases distratoras
        obtidas por semelhança fonética;
  \item avaliar a exatidão da ativação com um grupo maior de pessoas, sotaques e níveis de ruído;
  \item detectar automaticamente os pontos do rosto em novas fotografias, dispensando o cadastro manual;
  \item executar o agendador como serviço do sistema, para anunciar lembretes mesmo com a janela fechada;
  \item sincronizar a memória entre computadores com criptografia de ponta a ponta, mantendo a cadeia de
        integridade;
  \item publicar o pacote em um repositório próprio, com assinatura, para atualizações automáticas pelo
        gerenciador de pacotes;
  \item adicionar novos idiomas, aproveitando o catálogo e o teste de completude das traduções.
\end{enumerate}


% =============================================================================
%  Elementos pós-textuais (NBR 14724): referências (NBR 6023:2018), glossário e
%  apêndices.
% =============================================================================

% ---------- Referências ----------
\titulopostextual{Referências}
\begingroup
\singlespacing
\setlength{\parindent}{0pt}
\setlength{\parskip}{12pt}
\raggedright

ALPHA CEPHEI. \textbf{Vosk speech recognition toolkit}. [\textit{S. l.}]: Alpha Cephei, 2026.
Disponível em: \url{https://alphacephei.com/vosk/}. Acesso em: 22 out. 2026.

ANTHROPIC. \textbf{Claude API documentation}. San Francisco: Anthropic, 2026. Disponível em:
\url{https://docs.anthropic.com/}. Acesso em: 22 out. 2026.

ASSOCIAÇÃO BRASILEIRA DE NORMAS TÉCNICAS. \textbf{NBR 6023}: informação e documentação: referências:
elaboração. Rio de Janeiro: ABNT, 2018.

ASSOCIAÇÃO BRASILEIRA DE NORMAS TÉCNICAS. \textbf{NBR 10520}: informação e documentação: citações em
documentos: apresentação. Rio de Janeiro: ABNT, 2023.

ASSOCIAÇÃO BRASILEIRA DE NORMAS TÉCNICAS. \textbf{NBR 14724}: informação e documentação: trabalhos
acadêmicos: apresentação. Rio de Janeiro: ABNT, 2011.

BECK, Kent. \textbf{Test-driven development}: by example. Boston: Addison-Wesley, 2003.

COCKBURN, Alistair. \textbf{Hexagonal architecture}. [\textit{S. l.}], 2005. Disponível em:
\url{https://alistair.cockburn.us/hexagonal-architecture/}. Acesso em: 22 out. 2026.

COHN, Mike. \textbf{Succeeding with agile}: software development using Scrum. Upper Saddle River:
Addison-Wesley, 2009.

DEBIAN. \textbf{Debian policy manual}. [\textit{S. l.}]: The Debian Project, 2026. Disponível em:
\url{https://www.debian.org/doc/debian-policy/}. Acesso em: 22 out. 2026.

EVANS, Eric. \textbf{Domain-driven design}: tackling complexity in the heart of software. Boston:
Addison-Wesley, 2004.

FORD, Neal; PARSONS, Rebecca; KUA, Patrick. \textbf{Building evolutionary architectures}: support
constant change. Sebastopol: O'Reilly Media, 2017.

FOWLER, Martin. \textbf{Inversion of control containers and the dependency injection pattern}.
[\textit{S. l.}], 2004. Disponível em: \url{https://martinfowler.com/articles/injection.html}. Acesso
em: 22 out. 2026.

FOWLER, Martin. \textbf{Refactoring}: improving the design of existing code. 2. ed. Boston:
Addison-Wesley, 2018.

FREEDESKTOP.ORG. \textbf{Desktop entry specification}: version 1.5. [\textit{S. l.}], 2020.
Disponível em: \url{https://specifications.freedesktop.org/desktop-entry-spec/latest/}. Acesso em:
22 out. 2026.

GAMMA, Erich \textit{et al.} \textbf{Design patterns}: elements of reusable object-oriented software.
Reading: Addison-Wesley, 1994.

HARRIS, Charles R. \textit{et al.} Array programming with NumPy. \textbf{Nature}, London, v. 585,
p. 357-362, 2020.

HUNT, Andrew; THOMAS, David. \textbf{The pragmatic programmer}: from journeyman to master. Reading:
Addison-Wesley, 1999.

INTERNATIONAL ORGANIZATION FOR STANDARDIZATION. \textbf{ISO/IEC 25010:2011}: systems and software
engineering: systems and software quality requirements and evaluation (SQuaRE): system and software
quality models. Geneva: ISO, 2011.

JURAFSKY, Daniel; MARTIN, James H. \textbf{Speech and language processing}: an introduction to natural
language processing, computational linguistics, and speech recognition. 2. ed. Upper Saddle River:
Prentice Hall, 2009.

LARMAN, Craig; BASILI, Victor R. Iterative and incremental development: a brief history.
\textbf{Computer}, Los Alamitos, v. 36, n. 6, p. 47-56, 2003.

LEVKIVSKYI, Ivan; LEHTOSALO, Jukka; LANGA, Łukasz. \textbf{PEP 544}: protocols: structural subtyping
(static duck typing). [\textit{S. l.}]: Python Software Foundation, 2017. Disponível em:
\url{https://peps.python.org/pep-0544/}. Acesso em: 22 out. 2026.

LISKOV, Barbara H.; WING, Jeannette M. A behavioral notion of subtyping. \textbf{ACM Transactions on
Programming Languages and Systems}, New York, v. 16, n. 6, p. 1811-1841, 1994.

MARTIN, Robert C. \textbf{Agile software development}: principles, patterns, and practices. Upper
Saddle River: Prentice Hall, 2003.

MARTIN, Robert C. \textbf{Código limpo}: habilidades práticas do agile software. Rio de Janeiro: Alta
Books, 2009.

MARTIN, Robert C. \textbf{Arquitetura limpa}: o guia do artesão para estrutura e design de software.
Rio de Janeiro: Alta Books, 2019.

McCABE, Thomas J. A complexity measure. \textbf{IEEE Transactions on Software Engineering}, New York,
v. SE-2, n. 4, p. 308-320, 1976.

MESZAROS, Gerard. \textbf{xUnit test patterns}: refactoring test code. Upper Saddle River:
Addison-Wesley, 2007.

MEYER, Bertrand. \textbf{Object-oriented software construction}. Englewood Cliffs: Prentice Hall,
1988.

NATIONAL INSTITUTE OF STANDARDS AND TECHNOLOGY. \textbf{FIPS PUB 180-4}: secure hash standard (SHS).
Gaithersburg: NIST, 2015.

NIELSEN, Jakob. \textbf{Usability engineering}. Boston: Academic Press, 1993.

OLLAMA. \textbf{Ollama API}. [\textit{S. l.}], 2026. Disponível em:
\url{https://github.com/ollama/ollama/blob/main/docs/api.md}. Acesso em: 26 out. 2026.

PIPEWIRE. \textbf{PipeWire documentation}. [\textit{S. l.}], 2026. Disponível em:
\url{https://docs.pipewire.org/}. Acesso em: 26 out. 2026.

POVEY, Daniel \textit{et al.} The Kaldi speech recognition toolkit. \textit{In}: IEEE WORKSHOP ON
AUTOMATIC SPEECH RECOGNITION AND UNDERSTANDING, 2011, Waikoloa. \textbf{Proceedings} [...].
[\textit{S. l.}]: IEEE Signal Processing Society, 2011.

PRESSMAN, Roger S.; MAXIM, Bruce R. \textbf{Engenharia de software}: uma abordagem profissional.
8. ed. Porto Alegre: AMGH, 2016.

PRESTON-WERNER, Tom. \textbf{Semantic versioning 2.0.0}. [\textit{S. l.}], 2013. Disponível em:
\url{https://semver.org/}. Acesso em: 22 out. 2026.

SOHN, Jongseo; KIM, Nam Soo; SUNG, Wonyong. A statistical model-based voice activity detection.
\textbf{IEEE Signal Processing Letters}, New York, v. 6, n. 1, p. 1-3, 1999.

SOMMERVILLE, Ian. \textbf{Engenharia de software}. 10. ed. São Paulo: Pearson Education do Brasil,
2018.

SQLITE. \textbf{SQLite documentation}. [\textit{S. l.}], 2026. Disponível em:
\url{https://www.sqlite.org/docs.html}. Acesso em: 22 out. 2026.

VASWANI, Ashish \textit{et al.} Attention is all you need. \textit{In}: CONFERENCE ON NEURAL
INFORMATION PROCESSING SYSTEMS, 31., 2017, Long Beach. \textbf{Advances in neural information
processing systems 30}. Red Hook: Curran Associates, 2017. p. 5998-6008.

\endgroup

% ---------- Glossário ----------
\titulopostextual{Glossário}
\begingroup
\singlespacing
\setlength{\parindent}{0pt}
\setlength{\parskip}{8pt}

\textbf{Adaptador}: componente que implementa uma porta usando uma tecnologia concreta, como um banco de
dados ou um serviço de voz.

\textbf{Avatar}: representação gráfica do assistente; neste trabalho, um busto gerado a partir de
fotografias do usuário.

\textbf{Cache LRU}: armazenamento temporário que, quando cheio, descarta primeiro os itens usados há
mais tempo.

\textbf{Cadeia de resumos}: sequência de registros em que cada um guarda o resumo criptográfico do
anterior, tornando detectável qualquer alteração.

\textbf{Dublê de teste}: objeto que substitui um colaborador real durante um teste, como um sintetizador
que apenas registra os textos recebidos.

\textbf{Fluxo contínuo} (\en{streaming}): entrega de um resultado em partes, à medida que é produzido.

\textbf{Frase de ativação}: expressão que chama o assistente sem uso das mãos; neste trabalho, “ok Fael”.

\textbf{Frase distratora}: frase parecida com a de ativação, incluída na gramática para absorver o que
não deve ativar o assistente.

\textbf{Função de aptidão arquitetural}: verificação objetiva, preferencialmente automatizada, de uma
característica desejada da arquitetura.

\textbf{Gramática restrita}: lista fechada de frases entre as quais o reconhecedor de fala deve escolher.

\textbf{Pipeline}: encadeamento de estágios que processam dados em sequência, de forma sobreposta no
tempo.

\textbf{Porta}: contrato que descreve o que a aplicação precisa do mundo externo, ou o que oferece a ele.

\textbf{Raiz de composição}: ponto único do programa em que as implementações concretas são criadas e
conectadas.

\textbf{Token de cancelamento}: sinal compartilhado pelos estágios de uma interação, cuja ativação
interrompe todos eles.

\textbf{\en{Wheel}}: formato de distribuição de pacotes Python já construídos, instaláveis sem
compilação.

\endgroup

% ---------- Apêndices ----------
\clearpage
\phantomsection
\renewcommand{\thequadro}{A.\arabic{quadro}}
\setcounter{quadro}{0}
\titulopostextual{Apêndice A – Rastreabilidade entre requisitos e testes}

Este apêndice associa cada requisito à suíte que o verifica. Os nomes das suítes correspondem aos
arquivos da pasta \cod{tests}.

\begin{singlespace}
\legendaquadro{Rastreabilidade dos requisitos funcionais}
\small
\begin{tabularx}{\textwidth}{|P{2.4cm}|L|}
\hline
\textbf{Requisitos} & \textbf{Suítes de teste} \\ \hline
\endfirsthead
\hline
\textbf{Requisitos} & \textbf{Suítes de teste} \\ \hline
\endhead
RF01, RF14 & \cod{test_avatar}, \cod{test_configuracoes} \\ \hline
RF02 & \cod{test_animacao}, \cod{test_analise_audio} \\ \hline
RF03, RF09, RF13 & \cod{test_orquestrador}, \cod{test_vad} \\ \hline
RF04, RF05 & \cod{test_ativacao}, \cod{test_ganho}, \cod{test_sessao_pipeline_vigia} \\ \hline
RF06 & \cod{test_adaptadores_voz}, \cod{test_perfil_voz}, \cod{test_configuracoes} \\ \hline
RF07, RF08, RF10, RF12 & \cod{test_bootstrap_ui}, \cod{test_orquestrador}, \cod{test_estados_barramento}, \cod{test_configuracoes} \\ \hline
RF11 & \cod{test_bootstrap_ui}, \cod{test_i18n} \\ \hline
RF15 a RF17, RF19 & \cod{test_agenda}, \cod{test_modelo_offline}, \cod{test_ferramentas} \\ \hline
RF18 & \cod{test_agendador}, \cod{test_modelo_offline} \\ \hline
RF20, RF22, RF23 & \cod{test_sqlite}, \cod{test_repositorio_json} \\ \hline
RF21 & \cod{test_memoria}, \cod{test_i18n} \\ \hline
RF24, RF26 & \cod{test_modelo_claude}, \cod{test_ferramentas} \\ \hline
RF25, RF27 & \cod{test_modelo_offline}, \cod{test_bootstrap_ui} \\ \hline
RF28, RF29 & \cod{test_programas} \\ \hline
RF30 & \cod{test_perfil_voz}, \cod{test_analise_audio} \\ \hline
RF31, RF32 & \cod{test_configuracoes}, \cod{test_sessao_pipeline_vigia} \\ \hline
RF33, RF34 & \cod{test_i18n} \\ \hline
RF35 & \cod{test_configuracao_tema}, \cod{test_i18n} \\ \hline
RF36 & roteiro \cod{empacotamento/testar_deb.sh} (\cod{make testar-deb}) \\ \hline
RF37 & \cod{test_ollama} \\ \hline
RF38 & \cod{test_banco_e_modelos}, \cod{test_configuracoes} \\ \hline
RF39, RF40 & \cod{test_microfone_inicializacao}, \cod{test_configuracoes} \\ \hline
RF41, RF42 & \cod{test_banco_e_modelos}, \cod{test_configuracoes} \\ \hline
\end{tabularx}
\fonte{elaboração própria (2026).}
\end{singlespace}

\begin{singlespace}
\legendaquadro{Rastreabilidade dos requisitos não funcionais}
\small
\begin{tabularx}{\textwidth}{|P{2.4cm}|L|}
\hline
\textbf{Requisitos} & \textbf{Verificação} \\ \hline
\endfirsthead
\hline
\textbf{Requisitos} & \textbf{Verificação} \\ \hline
\endhead
RNF01 & \cod{test_bootstrap_ui} e medição de tempo \\ \hline
RNF02, RNF03 & \cod{test_orquestrador}, \cod{test_sessao_pipeline_vigia} \\ \hline
RNF04 & \cod{test_agendador} \\ \hline
RNF05 & \cod{test_animacao}, \cod{test_avatar}, \cod{test_adaptadores_voz} \\ \hline
RNF06 & \cod{test_sqlite} \\ \hline
RNF07 & \cod{test_modelo_claude} \\ \hline
RNF08 & \cod{test_bootstrap_ui}, \cod{test_ativacao}, \cod{test_adaptadores_voz} \\ \hline
RNF09 & \cod{test_orquestrador}, \cod{test_estados_barramento} \\ \hline
RNF10, RNF12, RNF17 & revisão de código e do catálogo de mensagens \\ \hline
RNF11 & \cod{test_ferramentas}, \cod{test_programas} \\ \hline
RNF13 & \cod{test_ativacao} \\ \hline
RNF14 a RNF16 & \cod{test_configuracoes}, \cod{test_bootstrap_ui} \\ \hline
RNF18 & \cod{test_i18n} \\ \hline
RNF19, RNF22 & \cod{test_arquitetura} \\ \hline
RNF20, RNF21 & ruff, mypy e pytest-cov em \cod{make verificar} e na integração contínua \\ \hline
RNF23, RNF24 & integração contínua (Python 3.10 e 3.12) e \cod{make testar-deb} \\ \hline
RNF25 & revisão do arquivo \cod{CHANGELOG.md} \\ \hline
RNF26, RNF27 & \cod{test_renderizacao}, \cod{test_analise_audio}, \cod{test_adaptadores_voz} e medição no app real \\ \hline
RNF28 & \cod{test_banco_e_modelos} e a fixture \cod{casa_isolada} (\cod{tests/conftest.py}) \\ \hline
\end{tabularx}
\fonte{elaboração própria (2026).}
\end{singlespace}

\clearpage
\renewcommand{\thequadro}{B.\arabic{quadro}}
\setcounter{quadro}{0}
\titulopostextual{Apêndice B – Comandos de voz}

O Quadro~B.1 apresenta exemplos de comandos aceitos. Os comandos são entendidos nos dois idiomas,
independentemente do idioma das respostas. Com o modelo de linguagem ativo, qualquer formulação em
linguagem natural é aceita; os exemplos abaixo também funcionam no modo offline, exceto quando indicado.

\begin{singlespace}
\legendaquadro{Exemplos de comandos de voz}
\small
\begin{tabularx}{\textwidth}{|P{3.4cm}|L|L|}
\hline
\textbf{Função} & \textbf{Português} & \textbf{Inglês} \\ \hline
\endfirsthead
\hline
\textbf{Função} & \textbf{Português} & \textbf{Inglês} \\ \hline
\endhead
Chamar o assistente & ok Fael & ok Fael \\ \hline
Chamar e comandar & ok Fael, anota tarefa comprar pão & ok Fael, add task buy bread \\ \hline
Criar tarefa & anota tarefa revisar o relatório & add task review the report \\ \hline
Listar tarefas & quais são minhas tarefas? & what are my tasks? \\ \hline
Concluir tarefa & terminei a tarefa 2 & complete task 2 \\ \hline
Lembrete relativo & me lembra daqui a 20 minutos de beber água & remind me in 20 minutes to drink water \\ \hline
Lembrete em data e hora (requer o modelo) & amanhã às 9 me lembra da reunião & tomorrow at 9 remind me of the meeting \\ \hline
Histórico & o que eu fiz ontem? o que concluí na semana passada? & what did I do yesterday? what did I do last week? \\ \hline
Abrir programa & abre o navegador; abrir a calculadora & open the browser; open the calculator \\ \hline
Hora e data & que horas são? que dia é hoje? & what time is it? what's the date? \\ \hline
Nova conversa & nova conversa & new conversation \\ \hline
Encerrar o modo contínuo & tchau & bye \\ \hline
\end{tabularx}
\fonte{elaboração própria (2026).}
\end{singlespace}

\clearpage
\renewcommand{\thequadro}{C.\arabic{quadro}}
\setcounter{quadro}{0}
\titulopostextual{Apêndice C – Configuração}

O Quadro~C.1 relaciona as variáveis de ambiente e as opções de linha de comando. A precedência é:
preferência salva na tela de configurações, depois opção de linha de comando ou variável de ambiente e,
por fim, valor padrão.

\begin{singlespace}
\legendaquadro{Variáveis de ambiente e opções de linha de comando}
\small
\begin{tabularx}{\textwidth}{|P{4.6cm}|P{3.2cm}|L|}
\hline
\textbf{Variável ou opção} & \textbf{Padrão} & \textbf{Efeito} \\ \hline
\endfirsthead
\hline
\textbf{Variável ou opção} & \textbf{Padrão} & \textbf{Efeito} \\ \hline
\endhead
\cod{ANTHROPIC_API_KEY} & -- & Ativa o Claude; sem ela, modelo local (Ollama) ou modo offline \\ \hline
\cod{ROBOFAEL_CEREBRO} & auto & \cod{auto} (Claude, depois Ollama, depois offline), \cod{claude}, \cod{ollama} ou \cod{offline} \\ \hline
\cod{ROBOFAEL_OLLAMA_URL} & \cod{http://localhost:11434} & Servidor do Ollama \\ \hline
\cod{ROBOFAEL_OLLAMA_MODELO} & \cod{llama3.2:3b} & Modelo local (sem ele, o primeiro instalado que aceite ferramentas) \\ \hline
\cod{ROBOFAEL_USUARIO} & Rafael & Nome usado na saudação e nas instruções do modelo \\ \hline
\cod{ROBOFAEL_ASSISTENTE} & Robofael & Nome do assistente \\ \hline
\cod{ROBOFAEL_FRASE_ATIVACAO} & ok Fael & Frase de ativação \\ \hline
\cod{ROBOFAEL_ATIVACAO} ou \cod{--sem-ativacao} & 1 (ligada) & Liga ou desliga a escuta passiva \\ \hline
\cod{ROBOFAEL_IDIOMA} ou \cod{--idioma} & pt-BR & Idioma das respostas e da interface (pt-BR ou en-US) \\ \hline
\cod{ROBOFAEL_MODELO} & \cod{claude-opus-5-5} & Modelo do Claude \\ \hline
\cod{ROBOFAEL_VOZ} & \cod{pt-BR-AntonioNeural} & Voz neural padrão \\ \hline
\cod{ROBOFAEL_DADOS} & \cod{./dados} (instalado: \cod{~/.local/share/robofael}) & Banco de dados (que pode ser movido nas configurações), perfil e gravações da voz \\ \hline
\cod{ROBOFAEL_CACHE} & \cod{./.cache} (instalado: \cod{~/.cache/robofael}) & Quadros do avatar e falas sintetizadas \\ \hline
\cod{ROBOFAEL_ASSETS} & \cod{./assets} & Fotografias dos avatares \\ \hline
\cod{ROBOFAEL_MODELOS} & \cod{./modelos} & Modelo de reconhecimento local \\ \hline
\cod{--avatar} & o salvo & Avatar inicial (crachá ou \en{selfie}) \\ \hline
\cod{--continuo} & desligado & Volta a ouvir após cada resposta \\ \hline
\cod{-v} & desligado & Registros detalhados (texto ouvido, pico e ganho) \\ \hline
\end{tabularx}
\fonte{elaboração própria (2026).}
{\noindent\small\singlespacing Nota: os nomes anteriores à versão 3.5.0 (\cod{EUROBO_*}) ainda são aceitos,
com aviso de depreciação, quando a variável \cod{ROBOFAEL_*} correspondente não está definida.\par}
\end{singlespace}

\clearpage
\titulopostextual{Apêndice D – Procedimentos de instalação, execução e verificação}

\textbf{Instalação pelo pacote (Ubuntu 24.04).} Na pasta do projeto, executar \cod{make deb} para gerar
o pacote em \cod{dist/}. Em seguida, instalá-lo com
\cod{sudo apt install ./dist/robofael_3.8.0_amd64.deb}. O gerenciador de pacotes instala as dependências
de sistema, e o pacote cria o ambiente Python em \cod{/opt/robofael} sem acessar a internet. O
assistente é aberto pelo item Robofael do menu de aplicativos ou pelo comando \cod{robofael}; a
ferramenta de gravação da voz, pelo item correspondente ou pelo comando \cod{robofael-gravar-voz}. Para
remover, usar \cod{sudo apt remove robofael}, que preserva os dados, ou \cod{sudo apt purge robofael},
que também apaga a pasta do programa; os dados pessoais só são removidos manualmente.

\textbf{Configurações de uso.} O botão de engrenagem abre a tela de configurações. Além de idioma,
avatar, voz e ativação por voz, ela permite escolher o microfone; marcar as opções de abrir o assistente
ao entrar no computador e de usar o comando \cod{robofael} em qualquer terminal (no pacote, o comando já
existe); ver o modelo de linguagem em uso e os modelos compatíveis do Claude e do Ollama; e, na seção de
memória, mudar o banco de pasta ou fazer um backup, operações que valem na hora. Para restaurar um
backup, basta fechar o assistente e copiar o arquivo do backup sobre o \cod{memoria.db}. Como o
assistente aberto pelo login não lê o \cod{~/.bashrc}, a chave do Claude, nesse caso, deve ser definida
em \cod{~/.config/environment.d/}.

\textbf{Execução em modo de desenvolvimento.} Executar \cod{make instalar} para criar o ambiente com as
dependências de desenvolvimento, \cod{make baixar-modelo-voz} para obter o modelo de reconhecimento local
e \cod{make executar} para abrir o assistente. Opções podem ser passadas com
\cod{make executar ARGS="--idioma en-US -v"}. A gravação da voz é feita com \cod{make gravar-voz}.

\textbf{Verificação.} O comando \cod{make verificar} executa a análise estática, a verificação de tipos
e os 643 testes; \cod{make testar-rapido} omite os testes de interface e de renderização;
\cod{make cobertura} gera o relatório de cobertura; \cod{make integridade} verifica a cadeia de resumos
do histórico; e \cod{make testar-deb} constrói o pacote e o testa em uma pasta temporária.

\textbf{Geração deste documento.} O texto desta monografia está em \cod{docs/monografia}, em LaTeX, e é
compilado com \cod{make monografia}, que produz o arquivo \cod{docs/monografia/monografia.pdf}.


\end{document}
